\documentclass[12pt,a4paper]{article}
\usepackage{geometry}
\geometry{
    a4paper,
    top=2.5cm, bottom=2.5cm,
    left=3cm,  right=3cm,
    marginparwidth=0pt
}
\usepackage{parskip}          % space between paragraphs, no indent
\usepackage{setspace}
\setstretch{1.15}
\usepackage{enumitem}
\setlist[enumerate]{leftmargin=2em, label=(\arabic*)}
\usepackage{booktabs}         % for case-law reference table
\usepackage{longtable}
\usepackage{array}
\usepackage{ragged2e}
\usepackage{fancyhdr}
\pagestyle{fancy}
\fancyhf{}
\renewcommand{\headrulewidth}{0.4pt}
\fancyhead[C]{\small\textit{\bn{পশ্চিমবঙ্গ ও অন্য রাজ্যের 277টি মামলার বিবরণ (2022 সালের দেওয়ানি আপিল নং 8238) 12ই সেপ্টেম্বর, 2023 [সঞ্জীব খান ও অরবিন্দ কুমার, উন্নয়ন] হেডনোটস-এর বিবরণঃ হাইকোর্টের আগে পশ্চিমবঙ্গ রাজ্য দ্বারা দায়ের করা আন্তঃ-আদালতের আবেদন খারিজ করে দেওয়া হয় এবং 2 নং যুগ্ম শিল্প সচিব, খনি মন্ত্রক, বা পশ্চিম বঙ্গের খনি ও খনিজ সম্পদ মন্ত্রকের যুগ্ম সচিবকে নির্দেশ দেওয়া হয়। আবেদনকারীদের অবস্থান-পশ্চিমবঙ্গ রাজ্য, যে তারা প্রশ্নবিদ্ধ জমির 20.87 একরের মালিক এবং এই প্রশ্নের উত্তরে, তাদের কাছে খনি ইজারা কার্যকর করার ক্ষেত্রে কোনও দ্বিধা নেই-এটি বর্ণিত অবস্থান হওয়ায়, যা আরও যুক্ত করা হয়েছে, উক্ত এলাকার জন্য খনি জায়গার অনুমোদনের ক্ষেত্রে কোনও পার্থক্য থাকা উচিত নয়।}}}
\fancyfoot[C]{\thepage}

%% Semantic commands
\newcommand{\sectionlabel}[1]{\bigskip\noindent\textit{\textbf{#1}}\par\smallskip}
\newcommand{\heldlabel}[1]{\noindent\textbf{#1}\par}
\newenvironment{judgequote}
  {\begin{quote}\small}
  {\end{quote}}
\newenvironment{citationblock}
  {\begin{quote}\small\itshape}
  {\end{quote}}
\newenvironment{caselawref}
  {\begin{quote}\small}
  {\end{quote}}

\usepackage{fontspec}
\usepackage{polyglossia}

\setdefaultlanguage{bengali}

\newfontfamily\bengalifont[
    Script=Bengali,
    Path = .,
    UprightFont = * ,
    BoldFont = * ,
    ItalicFont = * ,
    BoldItalicFont = *
]{Noto Serif Bengali.ttf}

\newcommand{\bn}[1]{{\bengalifont #1}}
\usepackage{fontspec}
\usepackage{polyglossia}

\setdefaultlanguage{bengali}

\newfontfamily\bengalifont[
    Script=Bengali,
    Path = .,
    UprightFont = * ,
    BoldFont = * ,
    ItalicFont = * ,
    BoldItalicFont = *
]{Noto Serif Bengali.ttf}

\newcommand{\bn}[1]{{\bengalifont #1}}
\begin{document}

\begin{center}
{\large\textbf{[2023] 12 S.C.R. 277 : 2023 INSC 824}}\\[4pt]
\bn{278 প্রবিধান) সংশোধনী আইন, 2015-সংশোধনের উদ্দেশ্য ও উদ্দেশ্য-তিন ধরনের পরিস্থিতিতে ব্যতিক্রম বা বিধিবদ্ধারাগুলির প্রয়োগঃ অনুমোদন/লক্ষ্যঃ 2015 সালের সংশোধনী আইনের উদ্দেশ্য এবং উদ্দেশ্য হল খনিজ সম্পদের বরাদ্দ নিলামের মাধ্যমে করা হয় তা নিশ্চিত করা-এই কারণেই এমএমডিআর আইন, 1957-এর 10-এ ধারার উপ-ধারা (1)-এ বলা হয়েছে যে,'12.01.2015'- এর আগে প্রাপ্ত সমস্ত আবেদন অযোগ্য হয়ে যাবে-এম. এম. ডি. আর আইন,1957-এর ধারা 10-ক-এর (2) উপ-ধারায় নির্দিষ্ট তিন ধরনের পরিস্থিতির ক্ষেত্রে ব্যতিক্রম অথবা বিধিবদ্ধ আবেদন প্রযোজ্য। এম. এম. ডি. আর আইন, 1957-এর 11-এ-র দ্বিতীয় বিভাগটি হল, যেখানে কোনও নজরদারি অনুমতিপত্র বা প্রত্যাশিত লাইসেন্স মঞ্জুর করা হয়েছে, সেই ক্ষেত্রে অনুমতিপত্রধারী বা লাইসেন্সধারীর কোনও কনি ইজারা দ্বারা প্রত্যাশিত ছাড়পত্র পাওয়ার অধিকার রয়েছে এবং রাজ্য সরকার সন্তুষ্ট হয় যে অনুমতিপত্র ধারক বা ছাড়পত্র গ্রহীতা 1957 সালের এমএমডিআর আইনের 10-এ ধারার উপ-ধারা (2)-এর (1) থেকে (4) পর্যন্ত উপ-ধারায় নির্দিষ্ট প্রয়োজনীয়তাগুলি মেনে চলেছেন। [অনুচ্ছেদ 14] খনি ও খনিজ (উন্নয়ন ও উন্নয়ন) আইন, 1957-খনি ও খনি (উন্নয়ন এবং উন্নয়ন) সংশোধনী আইন, 2015-পশ্চিমবঙ্গের বাণিজ্য ও শিল্প দপ্তরের উপ-সচিব ডলোমাইট খনির জন্য 76 একর জমির ক্ষেত্রে 1 নং-সর্বশ্রী চিরঞ্জীলাল (খনিজ) শিল্পের পক্ষে নির্দিষ্ট শর্ত সাপেক্ষে-কেন্দ্রীয় সরকারের অনুমোদনের প্রয়োজন আছে কি নাঃ\\[4pt]
v.
\bn{279 আনুষঙ্গিক/আদেশঃ যদিও বর্তমান ক্ষেত্রে, বিজ্ঞপ্তি নং এস. ও. 423 (ই) তারিখের পরে, ডলোমাইটকে একটি ছোট খনিজ হিসাবে বিজ্ঞপ্তি দেওয়া হয়েছিল এবং তাই, কেন্দ্রীয় সরকারের অনুমোদনের প্রয়োজন ছিল না কারণ'16.07.2015'তারিখের'মঞ্জুড়ীদেশ'- কে পূর্বশর্তের সাথে হেজ করা হয়েছিল, যার মধ্যে জমির মালিকদের প্রশ্নযুক্ত চিঠি (রায়ত) জমা দেওয়ার প্রয়োজনীয়তা অন্তর্ভুক্ত ছিল। খনির জন্য নয়, রায়ত জমি চাষাবাদ ইত্যাদির জন্য ব্যবহার করতে হবে-একবার খনির কাজ শুরু হয়ে গেলে রায়তরা জমি ব্যবহার করতে পারবে না-ডাব্লু. বি. এল. আর আইন, 1955-এর ধারা 2-এর উপ-ধারা (10)-এর পরিপ্রেক্ষিতে, কোনও রায়ত মানে যে কোনও উদ্দেশ্যে জমি ধারণকারী কোনও ব্যক্তি বা প্রতিষ্ঠান-তবে, জমির ক্ষেত্রে রায়তের অধিকার ডাব্লুবিএলআর আইনের 4 ধারার উপ-ধারার (2এ) পরিপ্রেক্ষিতে অন্য কোনও ব্যক্তিকে তার হোল্ডিং থেকে বালি খনন, খনন বা ব্যবহার করার অনুমতি দেয় না, বা কোনও ব্যক্তিকে ইট বা টাইলস তৈরির জন্য তার হোল্ডিংয়ের মাটি বা কাদামাটি খনন বা ব্যবহারের অনুমতি দেয়, রাজ্য সরকারের লিখিত পূর্ব অনুমতি ব্যতীত-শর্ত লঙ্ঘনের ক্ষেত্রে, নির্ধারিত কর্তৃপক্ষ রায়কে নোটিশ দেওয়ার পরে একটি সুযোগ দেখাতে পারে এবং জরিমানা করতে পারে। ডব্লিউ. বি. এল. আর আইন, 1955-এর ধারা 4-বি-তে বলা হয়েছে যে, যে কোনও জমির অধিকারী প্রত্যেক রায়িয়াত এমনভাবে জমি রক্ষণাবেক্ষণ ও সংরক্ষণ করবে যাতে এলাকা হ্রাস না হয় বা তার চরিত্র পরিবর্তন না করা হয় বা যে উদ্দেশ্যে জমি বসানো হয়েছিল বা পূর্বে অনুমোদন/আদেশ ব্যতীত অন্য কোনও উদ্দেশ্যে জমি রূপান্তরিত না হয়, যদি না কালেক্টরের লিখিত পূর্ব অনুমতি থাকে। [অনুচ্ছেদ 17] খনি ও খনিজ (উন্নয়ন ও উন্নয়ন) আইন, 1957-পশ্চিমবঙ্গ ভূমি সংস্কার আইন, 1956-ধারা 14 ওয়াই-এ মুঞ্জুরি আদি স্টেট অফ ওয়েস্ট বেঙ্গল বনাম সরবিশ্বরী চিরঞ্জিলা (মিনরা) ইন্দ্রি বাগড়ীহ।\\[4pt]
\bn{পশ্চিমবঙ্গের বাণিজ্য ও শিল্প বিভাগের উপ-সচিব ডলোমাইট খনির জন্য 76 একর জমির ক্ষেত্রে বাগান্দির 1 নং-সর্বশ্রী চিরঞ্জীলাল (খনিজ) ইন্ডাস্ট্রিজের পক্ষে 280 তারিখের'16.07.2015'জারি করেছিলেন। তিনি বলেন, জমির মালিকদের কাছ থেকে প্রয়োজনীয় চিঠি জমা দেওয়ার প্রয়োজনীয়তা সহ নির্দিষ্ট কিছু শর্তে (রায়িয়াত)-সেখানে উল্লিখিত অন্য একটি শর্ত ছিল অনুমতি নেওয়ার প্রয়োজনীয়তা এবং 1955 সালের মেমো আইনের 4 নং ধারার ভিত্তিতে সিদ্ধান্ত নেওয়া যায় না যে, প্রয়োজনীয় জমি ধরে রাখার জন্য ডাব্লু. বি. এল. আর আইন, 1955-এর 14-ওয়াই এবং যথাযথ কর্তৃপক্ষ থেকে জমির প্লটের জন্য রূপান্তর শংসাপত্র সরবরাহ করা হবে। উপ-জেলা ভূমি ও ভূমি সংস্কার দফতর, পুরুলিয়া কর্তৃক জারি করা 06.03.2017 তারিখের ভি/তাতার অধিকার/775/15-এ বলা হয়েছে যে, রাজস্ব নথি অনুসারে জমিটি'ডুঙ্গরি'হিসাবে নথিভুক্ত করা হয়েছিল-এর কারণ হল যে রায়িয়াত জমি খনির জন্য নয়-সুতরাং, একটি দ্বন্দ্ব হল যে, রায়িয়াত জমির অনুদান এবং একই জমির শ্রেণিবদ্ধকরণ হিসাবে'ডুংগরি'হল প্রথম সরকার-পরবর্তী সরকার, সংশ্লিষ্ট জমির মালিকদের (রায়িয়াত) দ্বারা প্রাপ্ত জমি সম্পর্কিত (06.03.2017) চিঠিটিকে শুধুমাত্র'1 নং আইন'- এর ভিত্তিতেই নথিভুক্ত করতে হবে, এবং 2016-র পরবর্তী আইনি বিষয়গুলির উপর সিদ্ধান্ত নিতে হবে যে, সেই আইনের সাথে সম্পর্কিত তথ্য পরিবর্তন করা হবে কি না। 1 সর্বজনীন আইন, 2016-এর 61 নম্বর বিধির বিধানের সুবিধা পাওয়ার অধিকারী-এটি যাচাই ও নিশ্চিত করা হয়নি-1998 সালে প্রথম পক্ষের নেতৃত্বে আবেদনটি এখনও ভাল থাকবে কিনা তা নিয়ে একটি সমস্যা দেখা দেবে, কারণ সেই সময়, যখন আবেদনটি দায়ের করা হয়েছিল, তখন কেন্দ্রীয় সরকারের অনুমোদনের প্রয়োজন ছিল-আরেকটি অসুবিধা হল যে ডব্লিউবিএমডিটিসিএলকে একটি পক্ষ হিসাবে অন্তর্ভুক্ত করা হয়নি, যদিও এটি সর্বদা প্রথম পক্ষের দ্বারা করা দাবির বিরোধিতা করছিল-তবে, এই বিষয়গুলি জারি করা নির্দেশের আলোকে পরীক্ষা করা হচ্ছে না-প্রত্যয়, যে স্পষ্ট আদেশটি এই সময়ে পাস করা যাবে না।\\[4pt]
\bn{281টি শহর ও অন্যান্য রেফারেন্সের তালিকা ভূষণ পাওয়ার অ্যান্ড স্টিলিমিটেড বনাম এস. এল. সিল, অতিরিক্ত সচিব এস. সি. আর. 149; ভূষণ পাওয়ার অ্যাণ্ড স্টিলিমিটেড এবং অন্যান্য বনাম উড়িষ্যা রাজ্য 240; ভূষণ পাওয়ার অ্যাণ্ড স্টীলিমিটেড বনাম রাজেশ ভার্মা, (2014) 5টি এসসিসি এসসিআর 411; রাজস্থান কো-অপারেটিভ ডেইরি ফেডারেশন লিমিটেড বনাম মহা লক্ষ্মী মিংগ্রেট মার্কেটিং সার্ভিস প্রাইভেট লিমিটেড ও অন্যান্য, (1996) 10 এসসিসি 7 এসসিসিআর 863-উলেখ করার জন্য রায়।\\[4pt]
\bn{2017 সালের 282 নং সিএএন নং 6596 সহ তথ্যগুলি বরং চেকার করা হয়েছে, যদিও বিস্তারিতভাবে লক্ষ্য করা প্রয়োজন। পশ্চিমবঙ্গের শিল্প, বাণিজ্য ও উদ্যোগ দপ্তরের যুগ্ম সচিব 2 নং এপিলকারি বা কোনও অনুমোদিত আধিকারিককে নির্দেশ দেওয়া হয়েছে যে 2 নং পক্ষের পক্ষে ইজারা খনির কাজ চালানোর জন্য-দীনেশ আগরওয়াল, একমাত্র মালিক-বাগনডিহ-এর 1 নম্বর সর্বশ্রী চিরঞ্জীলাল (খনিজ) ইন্ডাস্ট্রিজ। 2। তথ্যগুলি চেকার, যদিও বিশদভাবে লক্ষ্য করার প্রয়োজন রয়েছে। 2 নং 07.08.1985-তে, পশ্চিমবঙ্গ খনিজ উন্নয়ন ও বাণিজ্য নিগম লিমিটেড 2 নং ইউনিটের পক্ষে দীর্ঘমেয়াদী খনির অনুদানের জন্য একটি আবেদনের নেতৃত্ব দিয়েছিল, এবং ডাব্লিউ. বি. এল. ডি.-তে খনি এবং খনি প্লটগুলির জন্য সহকারী শাখা, ডল্লম. এম. সি. ডি-র পক্ষে, এবং 2 নং আই. ডি1-এর পক্ষেও অর্ডার জারি করা হয়েছিল। 1-বাগান্দিহর শ্রী চিরঞ্জীলাল (খনিজ) ইন্ডাস্ট্রিজ লিমিটেড পশ্চিমবঙ্গের খনি ও খনিজ অধিদপ্তরের পুরুলিয়া জোনের মাইনিং অফিসের ভারপ্রাপ্ত আধিকারিকের কাছে 76 একর জমিতে মুজা-খরিদুয়ারা, কুমারী এবং বোচে ডলোমাইট উত্তোলনের উদ্দেশ্যে একটি খনি ইজারা দেওয়ার জন্য 2001 সালেরিট জারি করে একটি আবেদন করেছিল। পশ্চিমবঙ্গের বাণিজ্য ও শিল্প বিভাগের যুগ্ম সচিব 26.03.2003 তারিখের আরেকটি আদেশে বলেন যে, খনির আবেদন 1 ফর শর্ট,'ডব্লিউবিএমডিটিসিএল'।\\[4pt]
\bn{বাগান্দিহর 1 নম্বর সরকারি চিরঞ্জীলাল (খনিজ) ইন্ডাস্ট্রিজের 283 নম্বরটি ডব্লিউ. বি. এম. ডি. টি. সি. এল-এর পূর্ববর্তী আবেদনে যে এলাকার জন্য আবেদন করা হয়েছিল, সেই এলাকার সঙ্গে বাগন্দিহর সরকারি চিরঞ্জিলাল (খনি) ইন্ডাস্ট্রিস 2003 সালের 7505 নম্বর রিটটি ওভারল্যাপ করে। যৌথ শিল্প সচিবের আদেশকে চ্যালেঞ্জ করে আবেদনটি প্রত্যাখ্যান করা হয়। 2 নম্বর সমষ্টিগতভাবে ক্ষতিগ্রস্ত, 1 নম্বর-সার্বভৌম চিরঞ্জিতলাল (মিনারেল) ইন্ডস্ট্রিস অব বাগান্দিহ-এর প্রতিনিধিরা 2003 সালে কলকাতা হাইকোর্টের যৌথ নির্দেশ নং 7505 (ডাব্লু) এবং বাণিজ্য বিভাগের সচিবের দ্বারা গৃহীত আদেশের শুনানিতে বলেছিলেন যে, পশ্চিম বঙ্গের শিল্প সচিব এবং পশ্চিমবঙ্গের বাণিজ্য সচিবের মধ্যে যৌথ আই. ডি/আই. 1 নম্বরটি পর্যালোচনা করা হয়েছিল। 1-বাগান্দিহর সর্বশ্রী চিরঞ্জীলাল (খনিজ) ইন্ডাস্ট্রিজ, এই বিষয়ে একমত হয়েছিল যে বাগনডিহর 1-নং সর্বেশিরা চিরঞ্জিলাল (খনি) শিল্পকে 76 একরের পুরো খনির এলাকা মঞ্জুর করা হবে এবং বাকি এলাকার জন্য ইজারাকে ডব্লিউবিএমডিটিসিএল-এর পক্ষে দেওয়া হবে। অন্য কোনও কারণ উল্লেখ করা হয়নি এবং উক্ত আদেশে নির্দেশিত করা হয়েছে। এইভাবে, 3 নং আইডি এবং 2 নং আইডি-র আবেদন প্রত্যাখ্যান করে 1 নং সর্বসিরা চিরঞ্জী লাল (খনিজ শিল্প)-এর আবেদন খারিজ করে দেওয়া হয়েছে। ফলস্বরূপ, বন মন্ত্রকের পক্ষ থেকে 76 একর জমির মাইনিং প্ল্যানের অনুমোদন, এবং 1 নং মিনিস্ট্রি অফ মাইনিং অফ ইন্ডিয়ার পক্ষ থেকে জারি করা বিভিন্ন আইডি, এবং মাইনিং ইন্ডাস্ট্রিজের 1 নম্বর আইডি, প্রত্যাহার করা হয়েছিল। 2. 5 যাইহোক, পশ্চিমবঙ্গের খনি শাখার বাণিজ্য ও শিল্প বিভাগের যুগ্ম সচিবের দ্বারা 13.10.2006 তারিখের আদেশ বাতিল বা প্রত্যাহার করা হয়েছিল, যাতে রেকর্ড করা হয়েছিল যে এই আদেশটি জমির সঠিক অবস্থানির্ধারণ না করেই এবং এই সত্যের অজ্ঞতার কারণে যে 13.03.2003 এবং 26.03.2003-এর প্রত্যাখ্যানের আদেশগুলি ইতিমধ্যে 2003 সালের 7505 (ডাব্লু) নম্বর রিট আদালতে উচ্চ আদালতে চ্যালেঞ্জ করা হয়েছিল। কর্তৃপক্ষ পশ্চিমবঙ্গ বনাম সার্বভৌম চিরঞ্জিলাল (খনিজ) শিল্পের অবস্থা নির্ধারণ করেনি।\\[4pt]
\end{center}
\bigskip\hrule\bigskip

\bigskip\hrule\bigskip


\noindent [2023] 12 S.C.R. 277 : 2023 INSC 824


\noindent CASE DETAILS


\begin{caselawref}
\bn{মামলার 284 নং দফা। বাদীর পক্ষ থেকে প্রত্যাহারের আদেশকে চ্যালেঞ্জ করা হয়নি। 2। 2003 সালেরিট নম্বর 7505 (ডাব্লু)-এর নিষ্পত্তি করার সময়ও হাইকোর্টের নজরে আনা হয়নি। হাইকোর্টের এই আদেশে বলা হয়েছে যে, বিকল্প নম্বর 1-এর নেতৃত্বাধীন সম্পূরক অ্যাফি ডেভিড ফাই-এর উপর নির্ভর করে, যা ছিল বাগনডিহ-এর সর্বশ্রী চিরঞ্জীলাল (খনিজ) ইন্ডাস্ট্রিজ, যার তারিখ ছিল'13.10.2006>'। পশ্চিমবঙ্গের বাণিজ্য ও শিল্প বিভাগের যুগ্ম সচিব কর্তৃক গৃহীত 09.07.2014 তারিখের আদেশ অনুসারে, বাগনডিহ-এর 1-নং সর্বশ্রী চিরঞ্জীলাল (খনিজ) ইন্ডাস্ট্রিজের নেতৃত্বে আবেদনটি প্রত্যাখ্যান করা হয়েছিল। ডব্লিউবিএমডিটিসিএল-এর নেতৃত্বে পূর্ববর্তী আবেদনের উপর নির্ভর করে এই আদেশে উল্লেখ করা হয়েছে যে, 13.03.2003 এবং 26.03.2003-এর দুটি প্রত্যাখ্যানের আদেশ যুগ্ম সচিব তাঁর তারিখের 13.10.2006 আদেশের মাধ্যমে প্রত্যাহার করে নিয়েছিলেন। এই আদেশটি এই সত্যকেও নির্দেশ করে যে, লৌহ আকরিক, ম্যাঙ্গানিজ প্লটের জন্য ডাব্লুবিএমডিটিসিএল-এর সাথে সম্পর্কিত 07.04.1986-র তারিখের আবেদনটি বাতিল করা হয়েছিল, এবং এর আগে ডাব্লু. বি. এম. ডি. এল.-এর পক্ষ থেকে একটি সাধারণ সময়-ভিত্তিক প্রশ্ন এবং লংস্টোনের দ্বারা প্রত্যাখ্যাত হয়েছিল কিনা তা নিশ্চিত করার জন্য, এবং ডাব্লিউবিএম. টি. সি.-র পক্ষ 09.07.2014 তারিখের আদেশটি ইঙ্গিত করে যে ডাব্লুবিএমডিটিসিএল-এর বিরুদ্ধে বাটিল করণ এবং প্রত্যাখ্যানের বাগান্দিহর 1 নং-সর্বশ্রী চিরঞ্জীলাল (খনিজ) শিল্পের সাথে কিছু সম্পর্ক ছিল এবং সম্ভবত বাগান্ডির 1 নম্বর-সর্বোশ্রী চিরঞ্জিলাল (খনি) ইন্ডাস্ট্রিজের পক্ষে 13.10.2006-এর আদেশটি ছিল। এটি <আই. ডি. 2> এর আদেশে দেওয়া কারণ থেকে প্রতিফলিত হয়, যেখানে বলা হয়েছে যে, তারিখের 13.10.2006 প্রত্যাহারের আদেশ বাতিল বা প্রত্যাহার করা হয়েছিল।
\end{caselawref}

\noindent HEADNOTES


\begin{judgequote}
\bn{285 03.12.2010, তারিখের 13.03.2003 এবং 26.03.2003 প্রত্যাখ্যানের আদেশগুলি এখনও বৈধ ছিল এবং ডাব্লুবিএমডিটিসিএল দ্বারা ডলোমাইট এবং চুনাপাথরের জন্য 07.08.1985 তারিখের কনি ইজারা-র আবেদন এখনও বহাল রয়েছে। এরপরে, 09.07.2014-এর আদেশে উল্লেখ করা হয়েছে খনি 2 11-এর ধারা 112-এর উপ-ধারা (2)-এ। (2) উপ-ধারা (1)-এর বিধানাবলী, যেখানে রাজ্য সরকার অফিসিয়াল গেজেটে পুনর্বিবেচনার অনুমতিপত্র বা সম্ভাবনার অনুমতিপত্র দেওয়ার জন্য এলাকাটি অবহিত করেনি বা প্রত্যাশিত লাইসেন্স বা পরিস্থিতি অনুযায়ী একটি অনুমোদনের অনুমতিপত্রের জন্য আবেদন করতে ব্যর্থ হয়নি এবং দুই বা ততোধিক ব্যক্তি, যা ক্ষেত্রমত, পুনর্বিবেচনা অনুমতিপত্র অথবা সম্ভাবনার লাইসেন্সের মেয়াদ শেষ হওয়ার তিন মাসের মধ্যে বা উল্লিখিত সরকার কর্তৃক বর্ধিত হওয়া অনুরূপ নির্দিষ্ট সময়ের মধ্যে, একটি অনুমতির জন্য আবেদন করেছে। এই শর্তে যে, যেখানে কোনও এলাকা, ক্ষেত্রমত, নজরদারি অনুমতিপত্র, প্রত্যাশিত লাইসেন্স বা হানি ইজারা মঞ্জুর করার জন্য উপলব্ধ রয়েছে এবং রাজ্য সরকার সরকারী গেজেটে এই ধরনের অনুমতিপত্র বা ইজারা প্রদানের জন্য বিজ্ঞপ্তি দ্বারা আবেদনপত্র আহ্বান করেছে, সেখানে নির্দিষ্ট সময়ের মধ্যে প্রাপ্ত সমস্ত আবেদনপত্র, যা এই ধরনের বিজ্ঞপ্তিতে উল্লেখ করা হয়েছে এবং যে আবেদনপত্রগুলি এই ধরনের বিজ্ঞপ্তি প্রকাশের আগে প্রাপ্ত হয়েছিল এবং সেগুলির নিষ্পত্তি করা হয়নি, সেগুলি এই উপধারার অধীনে অগ্রাধিকার দেওয়ার উদ্দেশ্যে একই দিনে প্রাপ্ত বলে মনে করা হবে। (3) উপ-ধারা (2)-এ উল্লিখিত বিষয়গুলি নিম্নরূপঃ-(ক) আবেদনকারীর দখলে থাকা পুনর্বিবেচনামূলক কাজকর্ম, সম্ভাব্য কাজকর্ম বা খনির কাজকর্ম সম্পর্কে যে কোনও বিশেষ জ্ঞান বা অভিজ্ঞতা; (খ) আবেদনের আর্থিক সম্পদ; (গ) আবেদনকারী কর্তৃক নিযুক্ত বা নিযুক্ত কর্মচারীদের প্রকৃতি ও গুণাগুণ; খনি ও খনিজ পদার্থের উপর ভিত্তি করে শিল্পে আবেদনকারী যে বিনিয়োগের প্রস্তাব দিয়েছেন তা উল্লেখ করা; (ঙ) নির্ধারিত অন্যান্য বিষয়।
\end{judgequote}

\begin{judgequote}
\bn{286 এবং খনিজ (উন্নয়ন ও উন্নয়ন) আইন, 19573-এ বলা হয়েছে যে, যে ক্ষেত্রে রাজ্য সরকার অফিসিয়াল গেজেটে নজরদারি অনুমতি, ইজারা খনির জন্য সম্ভাব্যতা লাইসেন্স এবং দুই বা ততোধিক ব্যক্তি অনুমতিপত্র, লাইসেন্স বা খননের জন্য আবেদন করেননি, সেই ক্ষেত্রে যার আবেদন আগে প্রাপ্ত হয়েছিল সেই ব্যক্তির উপর অনুমতি, ছাড়পত্র বা ইজারা দেওয়ার অধিকার থাকবে, যার আবেদন পরে প্রাপ্ত হয়েছিল। আদেশে বলা হয় যে ডাব্লুবিএমডিটিসিএল এই অঞ্চলে ডলোমাইট এবং চুনাপাথর খননে খুব আগ্রহী এবং উক্ত তথ্যটি 09.07.2014.81-এর মাধ্যমে লিখিত চিঠিতে স্বীকার করেছে। হাইকোর্টের সামনে 2135 নং যুক্তরাষ্ট্রীয় বাণিজ্য বিভাগের সচিব চিরঞ্জীলালের দ্বারা জারি করা যৌথ আদেশকে চ্যালেঞ্জানিয়ে 2014 সালের যৌথ শিল্প দফতরের জারি করা আদেশ। এই নির্দেশটি 10.09.2014 তারিখের আদেশের মাধ্যমে নিষ্পত্তি করা হয়েছিল যে, যুগ্ম সচিব, যিনি 09.07.2014-এর আদেশটি পাস করেছিলেন, তিনি তাঁর উপর ন্যস্ত এখতিয়ার প্রয়োগ করতে ব্যর্থ হয়েছিলেন কারণ ডব্লিউবিএমডিটিসিএল-এর নেতৃত্বে দায়ের করা আবেদনগুলি তারিখের সাধারণ আদেশ অনুসারে প্রত্যাখ্যান করা হয়েছিল এবং তাই সেগুলি বিকল্প ছিল না। উচ্চ আদালত দ্বারা 76 একর জমির ক্ষেত্রে দীর্ঘমেয়াদী ইজারা দেওয়ার জন্য নির্দেশ জারি করা হয়েছিল। এখানে এটা প্রাসঙ্গিক হতে পারে যে, এই আদেশে লিপিবদ্ধ করা হয়েছে যে, (4) উপ-ধারা (1)-এর বিধানাবলী সম্পর্কিত ক্ষেত্রে, যেখানে রাজ্য সরকার অফিসিয়াল গেজেটে, ক্ষেত্রমত, নজরদারি অনুমতি, সম্ভাব্যতা লাইসেন্স বা খনির ইজারা মঞ্জুর করার জন্য বিজ্ঞপ্তি দেয়, সেই সময়কালে প্রাপ্ত সমস্ত আবেদন, যা ত্রিশ দিনের কম হবে না, একই সঙ্গে দ্বি-পক্ষীয় হবে, যেন একই দিনে এই জাতীয় সমস্ত আবেদন প্রাপ্ত হয়েছে এবং রাজ্য সরকার, উপ-ধারায় নির্দিষ্ট বিষয়গুলি গ্রহণ করার পরে, (3) অনুমোদনের অনুমতি প্রদান করতে পারে। (5) উপ-ধারা (2)-এ যা কিছুই থাকুক না কেন, কিন্তু উপ-দফা (1)-এর বিধানগুলি সংশোধন করে রাজ্য সরকার, যে কোনও বিশেষ কারণে নথিভুক্ত করার জন্য, এমন কোনও আবেদনকারীকে, যার আবেদন পরে প্রাপ্ত হয়েছিল, আগে প্রাপ্ত আবেদনকারীর চেয়ে অগ্রাধিকার দিয়ে একটি পুনর্বিবেচনার অনুমতি, সম্ভাব্য লাইসেন্স বা একটি কনি ইজারা, ক্ষেত্রমত, প্রদান করতে পারেঃ প্রথম তফসিলে নির্দিষ্ট খনিজ সম্পর্কিত বিষয়ে, এই উপ-ধারার অধীনে কোনও আদেশ পাস করার আগে কেন্দ্রীয় সরকারের পূর্ব অনুমোদন নেওয়া হবে।
\end{judgequote}

\begin{judgequote}
\bn{ডব্লিউ. বি. এম. ডি. টি. সি. এল-এর 287টি আবেদন অপ্রাপ্য ছিল। তবে, উক্ত রিট নির্দেশের পক্ষভুক্ত না করে প্রত্যাখ্যানের আদেশগুলি প্রত্যাহার করা হয়েছিল। উল্লেখযোগ্যভাবে, ডব্লিউবিএমডিটিসিএল-এর নেতৃত্বে আবেদনটি, প্রযোজ্য নিয়মাবলির পরিপ্রেক্ষিতে সময়ের পরিপ্রেক্ষিতে অগ্রাধিকার দেওয়া হয়েছিল, যেখানে আবেদনটি প্রতিনিধি নম্বর 1-সর্বশ্রী চিরঞ্জিলাল (খনিজ) ইন্ডাস্ট্রিজ অফ বাগান্দিহর দ্বারা পরিচালিত হয়েছিল। 423 (ই), ডলোমাইটকে একটি ক্ষুদ্র খনিজ হিসাবে ঘোষণা করা হয়েছিল এবং সেই অনুযায়ী এখন থেকে এটি রাজ্য সরকারের আইনসভা ও প্রশাসনিক এখতিয়ারের অধীনে পড়ে। 2.10। পশ্চিমবঙ্গের বাণিজ্য ও শিল্প বিভাগের উপ-সচিব ডোলোমাইট খনির জন্য 1 নং-সর্বশ্রী চিরঞ্জিলাল (খনিজ) ইন্ডাস্ট্রিজের পক্ষে 76 একর জমির ক্ষেত্রে, 14-এর অধীনে উল্লিখিত অনুমতির জন্য পশ্চিমের 14-4 ধারার খসড়া অন্তর্ভুক্ত করার শর্তে একটি আই. ডি. 1 জারি করেছিলেন। এই অধ্যায়ের বিধানাবলীর ধারা 1-এর পরে যে কোনও সময়ে কোনও রায়াতের মালিকানাধীন জমির মোট আয়তন ধারা 14-এম-এর অধীনে প্রযোজ্য সর্বোচ্চ আয়তনের চেয়ে বেশি হলে, হস্তান্তর, উত্তরাধিকার বা অন্য কোনও কারণে, সেই জমির যে আয়তন ছাদ এলাকার বেশি হবে তা রাজ্যে নিহিত থাকবে এবং ছাদ এলাকা সম্পর্কিত এই অধ্যায়ে সমস্ত বিধান সেই জমির ক্ষেত্রে প্রযোজ্য হবেঃ যে কোনও ব্যক্তি চা বাগান, কল, কারখানা বা ওয়ার্কশপ, গবাদি পশু প্রজনন খামার, হাঁস-মুরগির খামার বা দুগ্ধ খামার অথবা জনপদ প্রতিষ্ঠা করার জন্য এই ধারার বিধিনিষেধ আরোপ করা হয়েছে।
\end{judgequote}

\begin{judgequote}
\bn{ডাব্লু. বি. এল. আর আইন, 1955-এর ধারা 4-সি 6-এর পরিপ্রেক্ষিতে যথাযথ কর্তৃপক্ষের কাছ থেকে জমির প্লটের জন্য প্রয়োজনীয় জমি দখল এবং রূপান্তর শংসাপত্র সরবরাহের জন্য 288 বঙ্গীয় ভূমি সংস্কার আইন, 5555। এটিও নির্ধারিত হয়েছে যে,'মুঞ্জুরি আদেশ'এবং'ইজারার দলিল'- এর পরবর্তী কার্যকরকরণ হল'নো অবজেকশন সার্টিফিকেশন'যা কেন্দ্রীয় সরকারের কাছ থেকে নেওয়া হবে কারণ উচ্চ আদালত কর্তৃক গৃহীত'আই. ডি. 1'তারিখের আদেশের সময় ডলোমাইট একটি প্রধান খনিজ ছিল। এই শর্তে বলা হয়েছে যে, রাজ্য সরকার কর্তৃক অনুমতিপ্রাপ্ত ব্যক্তি যদি অনুমতির তারিখের দুই বছরের মধ্যে সেই উদ্দেশ্যে জমি অধিগ্রহণ ও ধারণের জন্য ব্যবহার না করেন, যার জন্য রাজ্য সরকার তাঁকে অনুমতি দিয়েছে, তা হলে এই অধ্যায়ের ছাদ এলাকা সম্পর্কিত সমস্ত বিধান সেই জমির ক্ষেত্রে প্রযোজ্য হবে যা ধারা 14-এম-এর অধীনে তাঁর জন্য প্রযোজ্য ছাদ এলাকার অতিরিক্ত। জমির আয়তন, চরিত্র বা ব্যবহারের অনুমতিঃ-(1) যে উদ্দেশ্যে জমি বসানো হয়েছিল বা আগে ব্যবহার করা হচ্ছিল, সেই উদ্দেশ্য ব্যতীত অন্যে কোনও উদ্দেশ্যে সেই জমির এলাকা বা চরিত্র পরিবর্তন বা রূপান্তরের জন্য বা সেই জমি ব্যবহারের পদ্ধতিতে পরিবর্তনের জন্য কোনও জমির অধিকারী কোনও রায়ত কালেক্টরের কাছে আবেদন করতে পারবেন। (2) এই ধরনের আবেদন প্রাপ্তির পর, কালেক্টর, নির্ধারিত তদন্ত করার পরে এবং আবেদনকারীকে বা কোনওভাবে আগ্রহী ব্যক্তিদের শুনানির সুযোগ দেওয়ার পরে, লিখিত আদেশের মাধ্যমে আবেদনটি প্রত্যাখ্যান করতে পারেন বা নির্ধারিত শর্তাবলীর ভিত্তিতে, ক্ষেত্রমত, সেই জমি পরিবর্তন, রূপান্তর বা পরিবর্তনের নির্দেশ দিতে পারেন। (3) উপ-ধারা (4) এর অধীনে প্রদত্ত প্রতিটি আদেশে, যে তারিখ থেকে সেই তারিখটি নির্দিষ্ট করা হবে, যা রূপান্তর বা পরিবর্তন হবে। (4) উপ-ধারা (2) এর অধীন পরিবর্তন, রূপান্তর বা পরিবর্তন, যদি থাকে, নির্দেশ করে কালেক্টর কর্তৃক প্রদত্ত আদেশের একটি অনুলিপি অথবা তা থেকে আপিলে তিনি, ক্ষেত্রমত, 50 বা 51 ধারায় উল্লিখিত রাজস্ব দপ্তরের কাছে প্রেরণ করবেন এবং ঐরূপ রাজস্ব দপ্তর, ঐরূপ আদেশের দ্বারা উদ্ভূত অধিকারের নথিতে পরিবর্তনগুলি অন্তর্ভুক্ত করবে এবং ঐ আদেশ অনুসারে অধিকারের নথি সংশোধন করবে। (5) যদি কালেক্টর সন্তুষ্ট হন যে, যে উদ্দেশ্যে জমি নিষ্পত্তি করা হয়েছিল বা পূর্বে নির্ধারিত ছিল, তা ব্যতীত অন্য কোনও উদ্দেশ্যে জমি আহ্বান করা হচ্ছে।
\end{judgequote}


% --- page 2 ---

\noindent 278 SUPREME COURT REPORTS [2023] 12 S.C.R.


\begin{judgequote}
\bn{2014 সালের 289,1358 (ডাব্লু)। এই অবমাননার পিটিশনগুলি নিষ্পত্তি করা হয়েছিল, পর্যবেক্ষণ করে যে বাগান্দিহর 1 নং-সর্বশ্রী চিরঞ্জীলাল (খনিজ) ইন্ডাস্ট্রিজকে ডাব্লুবিএলআর আইন, 1955-এর ধারা 4-সি-এর অধীনে রূপান্তর শংসাপত্র সরবরাহ সহ শর্তগুলি পূরণ করতে হবে এবং ভারত সরকার থেকে অনাপত্তি শংসপত্র। অতএব, আদালত খুঁজে পেয়েছে যে 10.09.2014 তারিখের আদেশের কোনও ইচ্ছাকৃত, বা আপত্তিকর লঙ্ঘন হয়নি। যাইহোক, পশ্চিমবঙ্গ শিল্পের পক্ষকারের কাছে এই রিটটি পছন্দ করা হয়নি। ইতিমধ্যে, পশ্চিমবঙ্গের বাণিজ্য ও শিল্প দপ্তরের উপ-সচিব একটি ব্যাখ্যা চেয়েছিলেন এবং খনি মন্ত্রকের ভারত সরকার কর্তৃক জারি করা'26.08.2016'তারিখের একটি স্পষ্টীকরণের মাধ্যমে এটি স্পষ্ট করা হয়েছিল যে,'10.02.2015'- এর আগেও ডলোমাইট একটি অ-তফসিলি প্রধান খনিজ ছিল, যার জন্য 1957 সালের এমএমডিআর আইনের উপ-ধারা (1) থেকে ধারা 5-এর অধীনে কেন্দ্রীয় সরকারের পূর্বানুমোদন প্রয়োজন ছিল না। 'ডুঙ্গরি'হিসাবে শ্রেণীবদ্ধ জমি শুধুমাত্র খনি ইজারার উদ্দেশ্যে ব্যবহার করা হয় এবং এইভাবে, ডাব্লুবিএলআর, আইন, 1955-এর ধারা 4-সি-এর অধীনে রূপান্তর শংসাপত্রের প্রয়োজন নেই। অতিরিক্ত জেলা ম্যাজিস্ট্রেট এবং পুরুলিয়ার জেলা ভূমি ও ভূমি সংস্কার আধিকারিক দ্বারা জারি করা 07.04.2016 তারিখের স্পষ্টীকরণটি উল্লেখ করে যে, বাগনডিহর 1 নং-সরকারি চিরঞ্জীলাল (খনিজ) ইন্ডাস্ট্রিগুলি ইন্দ্রাণী রায় এবং রাজ্য সরকারের মালিকদের কাছ থেকে জমির একটি বড় অংশের বিষয়ে কোনও আপত্তি প্রমাণপত্র সংগ্রহ করেনি।
\end{judgequote}

\begin{judgequote}
\bn{290 <আই. ডি. 1> একর জমি, এইভাবে ডব্লিউ. বি. এল. আর আইন, 1955-এর ধারা 14-ওয়াই প্রযোজ্য হবে না কারণ বাগনডিহর সর্বশ্রী চিরঞ্জীলাল (খনিজ) ইন্ডাস্ট্রিজ ডব্লিউবিএলআর আইনের 14-এম ধারার অধীনে নির্ধারিত অতিরিক্ত সীমার মধ্যে জমি অধিগ্রহণ করেনি। 2014 সালের 21358 (ডাব্লু) সুবিদ ন্যাবলী, 2016 তারিখেরায়ের আগে ছিল। 3। আমরা পশ্চিমবঙ্গ রাজ্যের পক্ষে উপস্থিত বিদ্বান প্রবীণ আইনজীবী এবং 2 নং পক্ষের আইনজীবী-দীনেশ আগরওয়ালকে শুনেছি, যিনি ব্যক্তিগতভাবে উপস্থিত হয়েছেন। তারা তাদের লিখিত জমাও জমা দিয়েছেন। 4। আমরা আমাদের আলোচনা শুরু করি প্রথমে সুবেদ ন্যাবলি, 2016-এর বিধি 61 উল্লেখ করে, যা এই 12ই ফেব্রুয়ারির অধীনে পড়েঃ "61.Decleration সুবিদ নিয়াবলীর অযোগ্যতার বিষয়ে ছোটখাটো আবেদনের মধ্যে বড়খাটো খনিজ সম্পদের আবেদনগুলি অন্তর্ভুক্ত রয়েছে। শর্ত থাকে যে, আবেদনকারীকে যদি একটি'মঞ্জুর আদেশ'বা'লেটার অফ ইন্টেন্ট'(এল. ও. আই) বা অন্য কোনও সরকারি আদেশ জারি করা হয় যার জন্য আবেদনকারীর অবস্থান পরিবর্তনের প্রয়োজন হয়, তবে তার'হানি ইজারা'আবেদনটি সমস্ত প্রয়োজনীয় শর্ত যথাযথভাবে মেনে চলার পরে'বিবেচনার যোগ্য'হতে পারে।
\end{judgequote}

\begin{judgequote}
\bn{291 5. খনি ও খনিজ (উন্নয়ন ও উন্নয়ন) সংশোধনী আইন, 20158-এর মাধ্যমে 10-এ ধারা অন্তর্ভুক্ত করে এমএমডিআর আইন, 1957-এর প্রায় অনুরূপ একটি সংশোধনী করা হয়েছিল, যা নিম্নরূপঃ 10-ক। বিদ্যমান সুবিদার ধারক এবং আবেদনকারীদের অধিকার। (খ) খনি ও খনিজ (উন্নয়ন ও উন্নয়ন) সংশোধনী আইন, 2015-এর পূর্ববর্তী অংশে, যেখানে কোনও খনিজ সম্পদের জন্য কোনও জমির ক্ষেত্রে পুনর্বিবেচনার অনুমতিপত্র বা প্রস্পেক্টিং লাইসেন্স মঞ্জুর করা হয়েছে, সেখানে অনুমতিপত্রধারী বা লাইসেন্সধারীর কোনও খনিজাত সম্পদের বিষয়ে, সেই জমিতে সেই খনিজ সম্পত্তির ক্ষেত্রে, কেন্দ্রীয় সরকারের নির্ধারিত শর্তাবলী লঙ্ঘন না করে, এবং এই আইনের অধীনে অনুমোদনের অধিকারী বা লাইসেন্স গ্রহীতার দ্বারা, (i) খনিজাত সামগ্রীর অস্তিত্ব প্রতিষ্ঠার জন্য, এই জাতীয় মাপকাঠিতে, একটি প্রত্যাশিত লাইসেন্স পাওয়ার অধিকার থাকবে। এবং (4), ক্ষেত্রমত, নজরদারি অনুমতিপত্র বা প্রত্যাশিত লাইসেন্সের মেয়াদ শেষ হওয়ার তিন মাসের মধ্যে অথবা রাজ্য সরকার কর্তৃক বর্ধিত ছয় মাসের অনধিক সময়ের মধ্যে, সম্ভাব্য লাইসেন্স বা খনির ইজারা মঞ্জুরের জন্য আবেদন করতে ব্যর্থ হয়নি; 8 সংক্ষেপে,'সংশোধনী আইন, 2015'। পশ্চিমবঙ্গ বনাম সর্বশ্রী চিরঞ্জীল (খনি) ভারতেরাজ্য [সঞ্জীব খান্না, জে।]
\end{judgequote}


% --- page 3 ---

\begin{judgequote}
\bn{292 (গ) যদি কেন্দ্রীয় সরকার খনি ও খনিজ (উন্নয়ন ও উন্নয়ন) সংশোধনী আইন, 2015-এর প্রভুর সামনে খনি ইজারা মঞ্জুর করার জন্য ধারা 5-এর উপ-ধারা (1)-এর অধীনে প্রয়োজনীয় পূর্ববর্তী অনুমোদন জানিয়ে থাকে, অথবা রাজ্য সরকার কর্তৃক (যে নামে অভিহিত) অভিপ্রায়পত্র জারি করা হয়, তাহলে খনি ইজরা মঞ্জুরের জন্য 2016-র প্রথম ধারার অধীনে রাজ্যগুলিকে কোনও পূর্ববর্তী অনুমোদনের অনুমতি দেওয়া হবে না। 423 (ই) তারিখের 12.02.2015-এর আবেদনটি সুবিধার বিধান পরিষদ, 20169-কে কার্যকর আইন দেওয়ার আগে প্রাপ্ত হয়েছিল, তার মেয়াদ নির্বিশেষে অযোগ্য হয়ে যাবে। অন্য কথায়, এই আবেদনগুলি বৈধ হবে না। বিধানটি ব্যতিক্রম করে এবং বলে যে যদি কোনও আবেদনকারী, যিনি'29.07.2016'- এর আগে আবেদন করেছিলেন, তাকে অনুদান আদেশ বা অভিপ্রায় পত্র জারি করা হয়, বা আবেদনকারীর অবস্থান পরিবর্তনের প্রয়োজন হয় এমন অন্য কোনও আদেশ দেওয়া হয়, তবে তার আবেদনটি সমস্ত প্রয়োজনীয় শর্ত যথাযথভাবে মেনে চলার পরে হতে পারে। প্রশ্নটি হল'হাইকোর্টের দ্বারা মামলার শর্তাবলী অন্তর্ভুক্ত করা হয়েছে কি না'এবং 2016 সালের আইনের 61 নম্বর অনুচ্ছেদে বলা হয়েছে, যা 2016 সালের আইন অনুযায়ী প্রযোজ্য হবে। 1998 সালের মার্চ মাসে রিট আবেদনকারীদের আবেদনের ক্ষেত্রে এই ধরনের সাম্প্রতিক নীতি বা পশ্চিমবঙ্গ ক্ষুদ্র খনিজ সম্পদ আইন, 2016-র বিধানগুলি প্রযোজ্য হতে পারে না এবং আরও বেশি, যেহেতু যুগ্ম সচিবের ইজারা মঞ্জুর করার আদেশ 13ই অক্টোবর, 9 তারিখে জারি করা হয়েছে।
\end{judgequote}

\begin{judgequote}
\bn{293 2006 এবং এই আদালতের 10ই সেপ্টেম্বর, 2014 তারিখের দীর্ঘমেয়াদী ইজারা মঞ্জুর করার নির্দেশটি এই ধরনের নীতির আগে এবং উক্ত আইনটি কার্যকর হওয়ার আগে। এটি প্রতীয়মান হয় যে পরিবেশগত দিক বিবেচনা করে রিট আবেদনকারীদের দ্বারা প্রয়োজনীয় খনির পরিকল্পনা জমা দেওয়া হয়েছে এবং আপীলকারী/রাজ্য এ বিষয়ে কোনও অভিযোগ উত্থাপন করেনি। "7। পূর্বোক্ত অনুচ্ছেদে উল্লিখিত নীতিটি পশ্চিমবঙ্গ রাজ্যের প্রধান সচিব কর্তৃক জারি করা 02.02.2018 তারিখের চিঠির পরিপ্রেক্ষিতে রয়েছে, যেখানে নির্দিষ্টভাবে বলা হয়েছে যে রূপান্তর শংসাপত্র পাওয়া একটি শর্তের জন্য বাধ্যতামূলক এবং প্রত্যাহারের আদেশটি বাতিল করা হয়নি। 2003 সালের 7505 (ডাব্লু) একতার নিষ্পত্তি করা হয়েছিল, এটি লক্ষ্য না করেই যে 13.10.2006 তারিখের আদেশটি প্রত্যাহার করা হয়েছে বা বাতিল করা হয়েছে, যদিও রায়ে নির্দেশ দেওয়া হয়েছিল যে ইজারার অনুদানের আবেদনটি আইন অনুসারে হবে এবং প্রতিদানের সময় প্রযোজ্য হবে। 03.12.2010 এর আদেশটি কখনই প্রতিপক্ষ নং 1-এর দ্বারা চ্যালেঞ্জ করা হয়নি-বাগনডিহ-এর শ্রী চিরঞ্জিলাল (খনিজ) ইন্ডাস্ট্রিজ এবং ফি ন্যাশনালিটি অর্জন করেছে। সর্বোপরি, ভূষণের মামলাটি প্রতিপক্ষ লিমিটেডের আইন এবং মিনিস্ট্রিনের আগে রাজ্য খনি এবং মাইন্সের একটি অতিরিক্ত নম্বর 50-এর মধ্যে হবে। আবেদনটি ছিল ওড়িশার সম্বলপুর জেলায় একটি ইস্পাত কারখানা স্থাপনের প্রস্তাবের পরিপ্রেক্ষিতে। শ্রী ভূষণ পাওয়ার অ্যান্ড স্টিলিমিটেডকে কনি ইজারা অনুদানের প্রত্যাখ্যানকে উচ্চ আদালতে একটি রিট নির্দেশে চ্যালেঞ্জ করা হয়েছিল, যা খারিজ হয়ে গিয়েছিল, তবে এই আদালতের সামনে উত্থাপিত আপিলটি ভূষণ পাওয়ার অ্যাণ্ড স্টিল 10 (2017) 2 এস. সি. সি 125-এ 14.03.2012 তারিখেরায় দ্বারা অনুমোদিত হয়েছিল। স্টেট অফ ওয়েস্ট বেঙ্গল বনাম সর্বশ্রী চিরঞ্জীল (খনিজ) ইন্ডাস্ট্রিজ অফ বাগনডিহ [সঞ্জীব খান, জে]
\end{judgequote}

\begin{judgequote}
\bn{294 লিমিটেড এবং অন্যান্য বনাম উড়িষ্যা রাজ্য এবং অন্য 11, নিম্নলিখিত নির্দেশ সহ রাজ্য সরকারের তারিখ 09.02.2016-এর আদেশ বাতিল করেঃ "41। উপরের আলোকে, হাইকোর্ট এই রায় দিতে ভুল করেছে যে এটি পূর্ববর্তী সমঝোতাপত্রের অস্তিত্বের সময় সরকারের সাথে একটি নতুন সমঝোতা স্মারক স্বাক্ষর করার জন্য আবেদনকারীদের আহ্বান জানিয়ে রাজ্য সরকারের সিদ্ধান্তে হস্তক্ষেপ করতে পারে না। যেহেতু রাজ্য সরকার ইতিমধ্যে খনিজ সুবিদ নিয়ামক, 1960-এর শিথিলতার জন্য অন্যদের পক্ষে বরাদ্দ করেছে, তার নিয়ম 59 (2)-এর অধীনে, রাজ্যের পক্ষ থেকে আপিলকারীদের উক্ত বিশেষাধিকার অস্বীকার করার জন্য কোনও আধার তৈরি করা হয়নি। তদনুসারে, আমরা আপিলটি বাতিল করার অনুমতি দিই এবং হাইকোর্টের আদেশ এবং ব্লক কোর্টেরায়টিও খারিজ করে দিই। তদনুসারে, আমরা উড়িষ্যা রাজ্যকে 3য় তারিখের সমঝোতাপত্রের পরিপ্রেক্ষিতে যথাযথ পদক্ষেপ নেওয়ার নির্দেশ দিচ্ছি, এবং লাপাঙ্গায় তাদের ইস্পাত কারখানায় আবেদনকারীদের প্রয়োজনীয়তা পূরণের জন্য পর্যাপ্ত লৌহ আকরিক মজুদ মঞ্জুর করার জন্য কেন্দ্রীয় সরকারের কাছে আবেদনকারীদের মামলার সুপারিশ করার পূর্ববর্তী প্রতিশ্রুতিও দিচ্ছি। "9। এরপরে ওড়িশা রাজ্য পাওয়ার ভূষণ অ্যান্ড স্টিলিমিটেড এবং অন্যান্য বনাম ওড়িশা রাজ্য এবং অন্য 12 (বিকল্প) মামলারায় পর্যালোচনার জন্য একটি আবেদনের নেতৃত্ব দেয়, যা 2য় দিনের আদেশের মাধ্যমে প্রত্যাখ্যান করা হয়েছিল। 10। 1ম তারিখেরায় অমান্য ও নিষ্ক্রিয়তার অভিযোগ এনে, মেসার্স ভূষণ পাওয়ার অ্যান্ড স্টীলিমিটেডের নেতৃত্বে একটি আদালত অবমাননার মামলা দায়ের করা হয়। 2013 সালে ওড়িশা রাজ্য সরকার এই রায় পুনর্বিবেচনার বিরোধিতা করে। এই অবস্থান এই আদালতের পক্ষে ছিল না এবং রাজ্য সরকারের আধিকারিকদের 14.03.2012 11 (2012) 4 এসসিসি 246.12 (2011) 4এসসিসি 246.13 (2010) 13 এসসিসি 1 তারিখেরায় অবমাননার দায়ে দোষী সাব্যস্ত করা হয়।
\end{judgequote}

\begin{judgequote}
\bn{ভূষণ পাওয়ার অ্যান্ড স্টিলিমিটেড বনাম রাজেশ ভার্মা 14 মামলার 295 নম্বর রায়ের মাধ্যমে এই পরিস্থিতিতে রাজ্য সরকারকে আরও একটি সুযোগ দেওয়া হয়েছিল কেন্দ্রীয় সরকারের কাছে প্রয়োজনীয় সুপারিশ পাঠানোর জন্য। এই রায়টি পর্যবেক্ষণ করে যে এই আদালত এই সত্যটি উপেক্ষা করতে পারে না যে পক্ষগুলির মধ্যে একটি রায় রয়েছে, যা চূড়ান্ত হয়ে গেছে। তদনুসারে, এই যুক্তি যে সান্দুর ম্যাঙ্গানিজ এবং আয়রন ওরিস লিমিটেডের এই আদালতেরায়টি (সার্বভৌম) তারিখেরায়ে প্রদত্ত নির্দেশগুলি বাতিল করবে না। 14.03.2012 এর অধীনে পঠিত 22.04.2014 তারিখেরায়ের প্রাসঙ্গিক পর্যবেক্ষণগুলি প্রত্যাখ্যান করা হয়েছিল। আমরা এই সত্যটির দৃষ্টি হারাতে পারি না যে, এই রায়গুলির মধ্যে এমন কিছু অংশ রয়েছে যা দাবি করা হয়েছিল যে এই জমিগুলিতে হস্তক্ষেপ করা আবেদনগুলি খুব নির্দিষ্ট সময়ের মধ্যে খারিজ হয়ে গেছে, এমনকি যখন অন্যান্য আবেদনকারীদের হস্তক্ষেপও উল্লেখ করা হয়েছিল তখন তাদের হস্তক্ষেপও খারিজ হয়ে গিয়েছিল। যাইহোক, এই একই কারণে, এই আদালত রাজ্য সরকারকে উভয় ক্ষেত্রেই ইজারা খনির জন্য আবেদনকারীদের মামলার সুপারিশ করার জন্য নির্দেশ জারি করে। রায়ে দেওয়া এই ধরনের সুনির্দিষ্ট এবং যুক্তিসঙ্গত নির্দেশের পরিপ্রেক্ষিতে যা বাস্তবতা অর্জন করেছে, কেবলমাত্র সান্দুর ম্যাঙ্গানিজ মামলায় এই আদালত কর্তৃক অন্য একটি রায় দেওয়া হয়েছে বলে, এটি 14-3-2012 তারিখেরায়ে অন্তর্ভুক্ত নির্দেশগুলি বাতিল করার ভিত্তি হতে পারে না। একবার যখন আমরা ধরে নিই যে রাজ্য সরকার যতদূর পর্যন্ত আই. ডি. 1 তারিখেরায়ে অন্তর্ভুক্ত নির্দেশনা বাস্তবায়িত করতে বাধ্য, তখন তা মেনে চলতে বাধ্য এবং অন্যান্য প্রাসঙ্গিক বিবেচনার সাথে এই জাতীয় বিষয়গুলি রাজ্য সরকারের সুপারিশের ভিত্তিতে সিদ্ধান্ত নেওয়ার সময় কেন্দ্রীয় সরকারের বুদ্ধিমত্তার উপর ছেড়ে দেওয়া যেতে পারে। 14 (2014) 5 এস. সি. সি 551. স্টেট অফ ওয়েস্ট বেঙ্গল বনাম সর্বশ্রী চিরঞ্জিলাল (খনিজ) বাগান্দীর শিল্প [সঞ্জীব খান্না, জে।]
\end{judgequote}

\begin{judgequote}
\bn{যাইহোক, আমরা কেন্দ্রীয় সরকারের কাছে প্রয়োজনীয় সুপারিশ প্রেরণের মাধ্যমে অবমাননা দূরীকরণের জন্য আবেদনকারীদের একটি চূড়ান্ত সুযোগ দিচ্ছি। কেন্দ্রীয় সরকারের নিজস্ব যোগ্যতার ভিত্তিতে এবং আইন অনুসারে এই সুপারিশগুলি বিবেচনা করা হবে। যদি এই আদেশের অনুলিপি প্রাপ্তির তারিখ থেকে এক মাসের মধ্যে সুপারিশ পাঠানো হয়, তবে আমরা কোনও গুরুতর ব্যবস্থা না নেওয়ার প্রস্তাব দিচ্ছি এবং উত্তরদাতারা/অবমাননাকারীরা এই অবমাননা থেকে অব্যাহতি পাবেন। তবে, যদি উত্তরদাতাদের উপরে উল্লিখিত পদ্ধতিতে শুদ্ধ না হয়, তা হলে আবেদনকারীদের যথাযথ আবেদন পেশ করে এই আদালতে এটি নির্দেশ করার জন্য উন্মুক্ত থাকবে এবং সেই ক্ষেত্রে অবমাননাকারীদের বিরুদ্ধে ব্যবস্থা নেওয়া হবে। কেন্দ্রীয় সরকার অবশ্য এই অবস্থানিয়েছিল যে, এম. এম. ডি. আর আইন, 1957-এর সংশোধনীগুলি বিবেচনা করে, 2015-র সংশোধনী আইন 10-এ প্রবর্তনের মাধ্যমে, মেসার্স ভূষণ পাওয়ার অ্যান্ড স্টিলিমিটেডের অনুরোধটি অবৈধ হয়ে যায়। উপরোক্ত অবস্থানের পরিপ্রেক্ষিতে, কেন্দ্রীয় সরকারাজ্য সরকারকে চিঠি লিখেছিল, যার একটি অনুলিপি মেসর্স ভূষণ পাওয়ার অ্যাণ্ড স্টিলিমিটেডকে পাঠানো হয়েছিল। 13.05.2015 তারিখের চিঠিতে কেন্দ্রীয় সরকার বলেছিল যে খনিজ সহায়তার জন্য পূর্বানুমোদন অনুসারে প্রস্তাবটি অযোগ্য ছিল এবং তাই, এটিকে বন্ধ হিসাবে বিবেচনা করা উচিত। তবে, রাজ্য সরকার নিশ্চিত করতে পারে যে প্রস্তাবটি 1957-র এমএমডিআর আইনের ধারা 10-এর অধীনে অযোগ্যতার হাত থেকে নিরাপদ কিনা এবং সেই অনুযায়ী কেন্দ্রীয় সরকার 1957 সালের আইডির মাধ্যমে একই ধরনের পদক্ষেপ নিতে পারে। এই যোগাযোগের ফলস্বরূপ, রাজ্য সরকার 09.07.2015 তারিখের চিঠির মাধ্যমে মেসার্স ভূষণ পাওয়ার অ্যান্ড স্টিলিমিটেডকে জানিয়েছিল যে মণি ইজারা অনুদানের জন্য তাদের আবেদনগুলি এমএমডিআর আইন, 1957-এর ধারা 10-এ-এর উপ-ধারা (1) অনুসারে অযোগ্য হয়ে পড়েছে।
\end{judgequote}

\begin{judgequote}
\bn{297 12. মেসার্স ভূষণ স্টিল অ্যান্ড পাওয়ার লিমিটেড 15 (পূর্ববর্তী) মামলায় এই আদালত এই যুক্তিটি সুনির্দিষ্টভাবে পরীক্ষা করে দেখেছে যে, উক্ত মামলার তথ্যে, মেসার্স ইস্পাত ও বিদ্যুৎ লিমিটেড দ্বারা উত্থাপিত বিতর্কের পরিপ্রেক্ষিতে এমএমডিআর আইন, 1957-এর ধারা 10-এ-এর উপ-ধারা (2)-এর নির্দেশ (সি) প্রয়োগ করা যেতে পারে কিনা যে, রাজ্য সরকার কর্তৃক'খনি ইজারা'মঞ্জুর করার জন্য'লেটার অফ ইন্টেন্ট'জারি করা হয়েছিল এবং তাই, তাদের আবেদন সুরক্ষিত রয়েছে। জমা দেওয়া হয়েছিল যে,'24.05.2014'তারিখের সুপারিশটি রাজ্য সরকার দ্বারা'যে নামেই অভিহিত হোক না কেন'দ্বারা একটি লেটার অব ইন্টেন্ট হিসাবে বিবেচিত হওয়া উচিত, কারণ এটি নির্দেশ করে যে, 1957 সালের'এম. এম. ডি. আর.'আইনের অধীনে রাজ্য সরকার যতদূর পর্যন্ত নতুন অনুমতি পত্র জারি করতে পারে, ততদূর পর্যন্তা অনুমোদিত হতে পারে। উপরোক্ত যুক্তিগুলি মেসার্স ভূষণ স্টিল অ্যান্ড পাওয়ার লিমিটেডের (পূর্ববর্তী) ক্ষেত্রে এই আদালতের পক্ষে ছিল না এবং এই আদালতের পূর্ববর্তী তারিখেরায় এবং 2015 সালের সংশোধনী আইন কার্যকর হওয়ার আগে যে সমস্ত আবেদন কার্যকর হয়েছিল, তাতে এই পর্যবেক্ষণ সহ যে রাজ্য সরকার নির্দেশগুলি প্রয়োগ করতে ব্যর্থ হয়েছিল। এই আদালত মেসর্স ভূষণ স্টিল অ্যাণ্ড পাওয়ার লিমিটেড (পূর্ববর্ত)-এ জমা দেওয়া আবেদনগুলি প্রত্যাখ্যান করে এবং নিম্নোক্ত শর্তাবলীর অধীনেঃ "17. নিঃসন্দেহে, ধারা 10-এ-এর উপ-ধারা (1) অনুসারে, সমস্ত আবেদনগুলি বলবৎ হওয়ার আগে, 2015 সালের মধ্যে একটি প্রধান খনিজ আইনে পরিণত হওয়ার অযোগ্য ছিল। তখনকার পদ্ধতি অনুসারে, রাজ্য সরকার যে কোনও আবেদনকারীর দ্বারা কেন্দ্রীয় সরকারের কাছে কনি ইজারা অনুদানের জন্য জমা দেওয়া আবেদনের সুপারিশ করতে পারে এবং কেন্দ্রীয় সরকারকে অনুমোদন মঞ্জুর বা প্রত্যাখ্যান করার ক্ষমতা দেওয়া হয়েছিল। সুতরাং, ইজারা মঞ্জুর করার জন্য কেন্দ্রীয় সরকারের কাছ থেকে "পূর্ব অনুমোদন" অপরিহার্য ছিল, যা ছাড়া রাজ্য সরকার 15 (2012) 4 এস. সি. সি 246 প্রবেশ করতে পারেনি। পশ্চিম বঙ্গ রাজ্য বনাম সর্বশ্রী চিরঞ্জীল (খনিজ) বাগানের শিল্প [সঞ্জীব খান্না, জে।]
\end{judgequote}

\begin{judgequote}
\bn{আবেদনকারীর সঙ্গে এই ধরনের যে কোনও ইজারা চুক্তিতে 298 নম্বর। এই আদালত সেন্টার ফর পাবলিক ইন্টারেস্ট লিটিগেশন বনাম ইউনিয়ন অফ ইন্ডিয়া [(2012) 3 এস. সি. সি 1] (সংক্ষেপে "সি. পি. আই. এল মামলা") এবং প্রাকৃতিক সম্পদ বরাদ্দের ক্ষেত্রে প্রদত্ত রায়ে এই পদ্ধতির ত্রুটিগুলি লক্ষ্য করেছে। পুনর্বার, 2012 সালের 1 নম্বর বিশেষ জমি [প্রাকৃতিক সম্পদ বরাদ্দ, 2012-এর 1 নং বিশেষ জমি, 2012) 10 এস. চি 1]। এই রায়গুলিতে, এই আদালত প্রকাশ করেছে যে প্রাকৃতিক সম্পদের বরাদ্দ সাধারণত নিলামের মাধ্যমে হওয়া উচিত। সি. পী. আই এই বিবৃতি থেকে স্পষ্ট যে, বিদ্যমান আইন খনিজ সম্পদের ক্ষেত্রে ছাড়ের অনুমতি দেয় না বলে বিভিন্ন ধরনের সংস্কারের অভিজ্ঞতা হয়েছে। এটি দেখা গেছে যে, খনিজ সম্পদ বরাদ্দের ক্ষেত্রে স্বচ্ছতার উন্নতি হবে। সরকার খনিজ সম্পদগুলির মূল্য বৃদ্ধি করবে এবং প্রক্রিয়াগত বিলম্বকে প্রশমিত করবে, যা খনি ক্ষেত্রের বিকাশের উপর বিরূপ্রভাব ফেলবে। 19। 2015 সালের সংশোধনী আইনের লক্ষ্য হল, (1) বিচক্ষণতা দূর করা, (2) খনিজ সম্পদ বণ্টনে স্বচ্ছতা বাড়ানো, (3) প্রক্রিয়া সরলীকরণ, (4) প্রশাসনের ক্ষেত্রে বিলম্ব দূর করা যাতে দ্রুত ও সর্বোত্তম উন্নয়ন সম্ভব হয়। এবং, নতুন ধারা 8-এ, 9-বি অন্তর্ভুক্ত করা,
\end{judgequote}


% --- page 4 ---

\noindent 280 SUPREME COURT REPORTS [2023] 12 S.C.R.


\begin{judgequote}
\bn{299 9সি, 10এ, 10সি, 11বি, 11সি, 12এ, 15এ, 17এ, 20এ, 30বি, 30সি এবং চতুর্থ তফসিল 21। এই সংশোধনীগুলি প্রচলিত হয়েছেঃ (i) নিলামই হবে বরাদ্দের একমাত্র পদ্ধতি; (ii) বর্তমান ইজারার মেয়াদ শেষ পুনর্নবীকরণের তারিখ থেকে বাড়িয়ে 31-3-2030 (বন্দী খনিগুলির ক্ষেত্রে) এবং 31-3-2020 (বণিক খনি শ্রমিকদের জন্য) অথবা যতক্ষণ না ইতিমধ্যে নবায়ন মঞ্জুর করা হয়েছে, যদি থাকে, অথবা এই ধরনের ইজারকে মঞ্জুর করার তারিখ থেকে 50 বছরের সময়কাল পর্যন্ত; (iii) সংশ্লিষ্ট খনির ক্রিয়াকলাপ দ্বারা ক্ষতিগ্রস্ত ব্যক্তিদের স্বার্থ রক্ষার জন্য জেলা খনিজ ফাউন্ডেশন প্রতিষ্ঠা; (iv) খনিজ তহবিল থেকে জাতীয় খনিজ তহবিলের অবদান স্থাপন করা, ইত্যাদি, যাতে কেন্দ্রীয় খনি মালিকদের কাছ থেকে লৌহ আকরিকের অনুসন্ধান ও অপসারণের ক্ষেত্রে আরও সহজতর প্রক্রিয়ার অনুমোদন পাওয়া যায়। (vi) অবৈধ খনন প্রতিরোধে প্রতি হেক্টর 5 লক্ষ টাকা পর্যন্ত উচ্চতর জরিমানা এবং 5 বছর পর্যন্ত কারাদণ্ডের বিধান সম্বলিত কঠোর শাস্তিমূলক বিধান প্রবর্তন করা; এবং (vii) এই আইনের অধীনে অবৈধ সম্পদের বিচারের জন্য বিশেষ আদালত স্থাপনের জন্য রাজ্য সরকারকে ক্ষমতা প্রদান করা। "14। সুতরাং, 2015 সালের সংশোধনী আইনের উদ্দেশ্য ও উদ্দেশ্য হল খনিজ সম্পদের বরাদ্দ নিলামের মাধ্যমে করা হয় তা নিশ্চিত করা। এই কারণেই 1957 সালের এম. এম. ডি. আর আইনের 10-এ ধারার উপ-ধারা (1) নির্দেশ দেয় যে,'আই. ডি1'- এর আগে প্রাপ্ত সমস্ত আবেদন অযোগ্য হয়ে যাবে। দ্বিতীয় বিভাগটি হল, যেখানে কোনও নজরদারি অনুমতিপত্র বা সম্ভাব্য লাইসেন্স মঞ্জুর করা হয়েছে, সেই ক্ষেত্রে অনুমতিপত্রধারী বা লাইসেন্সধারীর কোনও খনি ইজারা দ্বারা সম্ভাব্য লাইসেন্স পাওয়ার অধিকার রয়েছে এবং রাজ্য সরকার এই মর্মে সন্তুষ্ট হয় যে, অনুমতিপত্রধারক বা ছাড়পত্রধারী এম. এম. ডি. আর আইন, 1957-এর উপ-ধারা (2) থেকে ধারা 10-এ-এর অংশ (খ)-এর (1) থেকে (4) পর্যন্ত উপ-ধারায় নির্দিষ্ট প্রয়োজনীয়তাগুলি মেনে চলেছেন। এই শ্রেণীর মামলাগুলি রক্ষা করার কারণ হল এই যে, তাঁরা অর্থ ব্যয় করে তাঁদের অবস্থান পরিবর্তন করেছিলেন।
\end{judgequote}

\begin{judgequote}
\bn{300। তদনুসারে, বৈধ প্রত্যাশার নীতি প্রয়োগ করা হয়। তৃতীয় বিভাগটি হল যেখানে কেন্দ্রীয় সরকার ইতিমধ্যে তাদের পূর্বের অনুমোদন জানিয়ে দিয়েছিল বা রাজ্য সরকার 2015 সালের সংশোধনী আইন কার্যকর হওয়ার আগে কনি ইজারা মঞ্জুর করার জন্য লেটার অফ ইন্টেন্ট জারি করেছিল। এর কারণ, এটি লক্ষ্য করা যায় যে এই আবেদনকারীদের কিছু অধিকার অর্জিত হয়েছিল কারণ সমস্ত প্রয়োজনীয় পদ্ধতি এবং আনুষ্ঠানিকতা মেনে চলা হয়েছিল এবং কেবল আনুষ্ঠানিক ইজারা কার্যকর করা বাকি ছিল। 15। রাজ্য সরকার কর্তৃক প্রদত্ত অনুমোদনের জন্য চিঠিটিকে যে কোনও নামের শব্দের উপর পূর্বাভাস দেওয়া অভিপ্রায় পত্র হিসাবে বিবেচনা করা যেতে পারে কিনা এই প্রশ্নের উত্তর দেওয়া। 'লেটার অফ ইন্টেন্ট'শব্দের অর্থের জন্য আইনি অভিধানে উললেখ করা হয়েছিল যে পক্ষগুলি একটি চুক্তি করতে চায় বা অন্য কোনও পদক্ষেপের জন্য একসাথে যোগ দিতে চায় তাদের মধ্যে প্রাথমিক বোঝাপড়া হিসাবে। এই আদালতের সিদ্ধান্তের ক্ষেত্রেও উলখ করা হয়েছিল ঋষি কিরণ লজিস্টিক্স প্রাইভেট লিমিটেড বনাম কান্ডলা পোর্ট্রাস্ট এবং অন্যান্যদের ট্রাস্টি বোর্ড 16 এবং রাজস্থান সমবায় ডেইরি ফেডারেশন লিমিটেড বনাম মহা লক্ষ্মী মিংগ্রেট মার্কেটিং সার্ভিস প্রাইভেট লিমিটেড এবং অন্যান্য 17 তবে, উক্ত যুক্তিটি প্রত্যাখ্যান করা হয়েছিলঃ "26। উপরোক্ত অর্থ প্রয়োগ করে, কি বলা যেতে পারে যে রাজ্য সরকারের 24-5-2014 তারিখের চিঠিটি একটি অভিপ্রায়ের চিঠি গঠন করবে? আমরা সহজ উত্তর দিতে ভয় পাচ্ছি, কারণটি নেতিবাচক, যেহেতু উপরে উল্লিখিত কোনও পূর্ববর্তী চুক্তিতে কেন্দ্রীয় সরকারের অনুমোদনের প্রয়োজনীয়তা উল্লেখ করা হয়েছিল। এই ধরনের অনুমোদন না পেলে, রাজ্য সরকার সম্ভাব্য লাইসেন্সধারী/ইজারা গ্রহীতার কাছে কোনও চুক্তি করার অভিপ্রায় জানাতে পারবে না কারণ পূর্ববর্তী অনুমোদনের অভাব থাকবে। অতএব, ভবিষ্যতে চুক্তি করার ইচ্ছা প্রকাশকারী কোনও আবেদনকারীকে রাজ্য সরকার কোনও প্রতিশ্রুতি দিতে পারবে না। পদটিতে 16 (2015) 13 এসসিসি 233. 17 (1996) 10 এসসিসি 405 থাকবে।
\end{judgequote}


% --- page 5 ---

\begin{judgequote}
\bn{301 পৃথক ছিল যদি কেন্দ্রীয় সরকারের পূর্ব অনুমোদন পাওয়ার পরে 24-5-2014 তারিখের চিঠি জারি করা হয়েছিল। তবে, তা নয়। কেন্দ্রীয় সরকারকে লেখা এই চিঠিটি কেবলমাত্র সুপারিশমূলক প্রকৃতির ছিল এবং চূড়ান্ত সিদ্ধান্তটি কেন্দ্রীয় সরকারের উপর ন্যস্ত ছিল। এটি একটি ভিন্ন বিষয় যদি কেন্দ্রীয় সরকার কোনও বহিরাগত কারণে তার অনুমোদন দিতে অস্বীকার করে বা দুর্বোধ্যতা বা আবেদন প্রত্যাখ্যান করার সময় প্রাসঙ্গিকারণ/উপাদান গ্রহণ না করে, যা একটি ভিন্ন মামলাগুলি তৈরি করতে পারে এবং আইনে উপলব্ধ ভিত্তিতে আবেদনটি প্রত্যাখ্যান করে কেন্দ্রীয় সরকারের পদক্ষেপকে চ্যালেঞ্জ করার কারণ হতে পারে। আবেদনকারীর আবেদনটি, অতএব, আইনের ধারা 10-এ-এর নির্দেশ (সি)-এর আওতাভুক্ত হবে না। 27। আমরা এই সত্যের পক্ষে সওয়াল করছি যে, আবেদনকারী প্রথমে এই আবেদনে সফল হয়েছিলেন, যেহেতু রাজ্য সরকারকে পক্ষগুলির মধ্যে স্বাক্ষরিত সমঝোতাপত্রের পরিপ্রেক্ষিতে আবেদনকারীর মামলাটি কেন্দ্রীয় সরকারের কাছে সুপারিশ করার জন্য নির্দেশ দেওয়া হয়েছিল। এটি করা হয়নি এবং সিদ্ধান্তটি ছিল ভূষণ পাওয়ার অ্যান্ড স্টিলিমিটেড বনাম রাজেশ ভার্মা-তে গৃহীত 22-4-2014 তারিখের আদেশে সম্পূর্ণরূপে অনুমোদিত। এটা সম্ভব যে রাজ্য সরকার দ্রুত পদক্ষেপ নিয়ে আগে সুপারিশগুলি পাঠিয়েছিল, কেন্দ্রীয় সরকার তার অনুমোদন দিতে পারত। তবে, এটি করা যেত কি না তা কেন্দ্রীয় সরকার কর্তৃক প্রদত্ত নির্দেশের বিপরীতে, স্টিলিমিটেডকে কখনও স্পষ্টভাবে পর্যবেক্ষণ করা হয়নি। রাজেশ ভার্মা বলেন যে, রাজ্য সরকারের সুপারিশগুলি তার নিজস্ব যোগ্যতার ভিত্তিতে এবং আইন অনুসারে বিবেচনা করা কেন্দ্রীয় সরকারের দায়িত্ব হবে। যদি কেন্দ্রীয় সরকার তা না করে থাকে, তবে এটি বর্তমান অবমাননার বিষয় হতে পারে না। "16। উপরোক্ত রায়টি আমাদের উদ্দেশ্যের জন্য প্রাসঙ্গিক, যদিও বর্তমান ক্ষেত্রে, 10.02.2015 তারিখের পোস্ট বিজ্ঞপ্তি নং. এস. ও. 423 (ই), ডলোমাইট স্টেট অফ ওয়েস্ট বেঙ্গল বনাম সর্বশ্রী চিরঞ্জীবী (খনিজ) বাগান্দির শিল্প [সঞ্জীব খান্না, জে.]
\end{judgequote}

\noindent LIST OF CITATIONS AND OTHER REFERENCES


\begin{citationblock}
\bn{কেন্দ্রীয় সরকারের অনুমোদনের প্রয়োজন ছিল না এই কারণে যে, ইজারার দলিল কার্যকর করার আগে জমির মালিকদের (রায়ত) সম্পত্তির চিঠি জমা দেওয়ার প্রয়োজনীয়তা সহ পূর্বশর্তের সাথে 16.07.2015 তারিখের মুঞ্জুরি দেশকে হেজ করা হয়েছিল, অথবা একটি শর্ত ছিল যে এই প্রভাবের একটি শর্ত খসড়া জারাতে অন্তর্ভুক্ত করা হবে। অতএব, আমাদের মতে, <আই. ডি. 1>-এর তারিখের মজুরি দেশটি অস্থায়ী, এবং এর শর্তাবলী পূরণ সাপেক্ষে। "" "(ক) এখন পর্যন্ত সংশোধিত (এম. সি. আইন, 1960-এর 1ম বিধি 31) হিসাবে খনিজ সুবিদায় নির্ধারিত মডেল ফর্মে আপনাকে একটি খসড়া নথি জমা দিতে হবে, (খ) খসড়াটি টেকসই কাগজে সুন্দরভাবে প্রস্তুত করতে হবে এবং দুটি লাইনের মধ্যে পর্যাপ্ত জায়গা রাখতে হবে, অনুমতি দেওয়ার জন্য, যদি প্রয়োজন হয়, সংশোধন করার জন্য, (গ) জারির পরে ইজারা দলিল আপনার নিজের খরচে নিবন্ধিত হবে এবং ন্যূনতম খনন কাজ শুরু করা উচিত নয়, আপনি 1974-এর 21 নম্বর খনির খনি প্রকল্প জমা করবেন, এবং যদি আপনি পরিবেশ সুরক্ষা আইন, 1986-এর ধারা 22 (4)-এর অধীনে অনুমোদিত মাইনিং প্ল্যান জমা না করে থাকেন, তা হলে আপনি নির্বাচন কমিশনের কাছে জমা দেবেন।" 10, 000/(দশ হাজার টাকা) শুধুমাত্র ইজারার শর্তাবলী যথাযথভাবে পালন করার জন্য সুরক্ষা হিসাবে,
\end{citationblock}

\begin{judgequote}
\bn{... 303 উপযুক্ত হিসাবের শীর্ষের অধীনে যা ইজারা চুক্তির মেয়াদ শেষ হওয়ার পরে আপনাকে ফেরত দেওয়া হবে, যদি না এর পুরো বা একটি অংশ আপনার পক্ষ থেকে কোনও আইনের জন্য সরকার দ্বারা আটকানো বা বাজেয়াপ্ত করা হয়, যার মধ্যে সরকারের প্রাপ্য অর্থ প্রদানের ক্ষেত্রে বরাদ্দ অন্তর্ভুক্ত থাকে [এম. সি. আইন, 1960-এর নিয়ম 32], (i) আপনাকে জমা দিতে হবে জমির মালিক (দের) যথাযথ চিঠি (গুলি)। নিয়ামবলি, 1960 সালের খসড়া ইজারা চুক্তিতে, (এম) দলিলটি কার্যকর করার আগে আপনাকে আপ টু ডেট রয়্যালটি ক্লিয়ারেন্স, ইনকাম ট্যাক্স ক্লিয়ার্যান্স এবং ভ্যাট ক্লিয়্যারেন্স শংসাপত্র জমা দিতে হবে, (এন) আপনাকে খসড়া দলিলের সাথে একটি জিও-উললেখ মানচিত্র জমা করতে হবে যা ডিএল অ্যান্ড এলআরও এবং ডিএমএম, পশ্চিমবঙ্গ দ্বারা যথাযথভাবে যাচাই করা হয়েছে, যদি জমা না দেওয়া হয়, (ও) আপনি প্রয়োজনীয় জমি ধরে রাখার জন্য ডাব্লুবিএলআর আইন, 1955-এর ধারা 14 ওয়াই-এর অধীনে অনুমতি পেয়েছেন, (পি) আপনাকে যথাযথ কর্তৃপক্ষ (ডাব্লু. বি. এল. আর আইনের ধারা 4 সি, 1955) থেকে জমির প্লটগুলির জন্য রূপান্তর শংসপত্র জমা নিতে হবে।
\end{judgequote}

\begin{judgequote}
\bn{304 (র) উপরোক্ত শর্তাবলী মেনে নির্ধারিত সময়ের মধ্যে দলিলটি কার্যকর না করা হলে ইজারা অনুমোদনের আদেশটি প্রত্যাহারের জন্য দায়বদ্ধ হবে, (গুলি) আপনাকে এই বিভাগের কাছে ইজারা কার্যকর করার দলিল পেশ করার আগে সমস্ত বিধিবদ্ধ প্রয়োজনীয়তা পূরণ করতে হবে। ডব্লিউ. বি. এল. আর আইন, 1955-এর ধারা 2-এর উপ-ধারা (10)-এর পরিপ্রেক্ষিতে, কোনও রায়িয়াত মানে যে কোনও উদ্দেশ্যে জমি ধারণকারী কোনও ব্যক্তি বা প্রতিষ্ঠান। তবে, ডব্লিউবিএলআর আইন, 1956-এর 4 ধারার উপ-ধারার (2এ) পরিপ্রেক্ষিতে জমির ক্ষেত্রে রায়িয়াতের অধিকার অন্য কোনও ব্যক্তিকে তার হোল্ডিং থেকে বালি খনন, খনন বা ব্যবহার করার অনুমতি দেয় না, বা কোনও ব্যক্তিকে রাজ্য সরকারের লিখিত পূর্ব অনুমতি ব্যতীত ইট বা টাইল তৈরির জন্য তার হোল্ডিংয়ের মাটি বা কাদামাটি খনন বা ব্যবহারের অনুমতি দেয়। শর্ত লঙ্ঘনের ক্ষেত্রে, নির্ধারিত কর্তৃপক্ষ কোনও রায়ায়াতকে নোটিশ দেওয়ার পরে এবং কারণ দেখানোর সুযোগ দেওয়ার পরে, আর্থিক জরিমানা আদায় করতে পারে। একটি আদেশ পাস হওয়ার পরে, জমি রাজ্য থেকে মুক্ত থাকবে। ডব্লিউ. বি. এল. আর আইন, 1955-এর ধারা 4-বি-তে বলা হয়েছে যে, যে কোনও জমির অধিকারী প্রত্যেক রায়িয়াত এমনভাবে জমি রক্ষণাবেক্ষণ ও সংরক্ষণ করবে যাতে এলাকা হ্রাস না হয় বা তার চরিত্র পরিবর্তন না করা হয় বা যে উদ্দেশ্যে জমি নিষ্পত্তি করা হয়েছিল বা পূর্বে অনুমোদিত/আদেশ ব্যতীত অন্য কোনও উদ্দেশ্যে জমি রূপান্তরিত না হয়, তবে কালেক্টরের লিখিত অনুমতি ব্যতীত। আমাদের উদ্দেশ্যের জন্য সমানভাবে তাৎপর্যপূর্ণ ডাব্লুবিএলআর আইনের ধারা 3 এ নয়, যা বলে যে সমস্ত অকৃষি ভাড়াটিয়া এবং আন্ডার-টেন্যান্টদের অধিকার এবং স্বার্থ রাজ্যে সমস্ত দায়বদ্ধতা থেকে মুক্ত থাকবে এবং পশ্চিমবঙ্গ এস্টেট অধিগ্রহণ আইন, 1953-র ধারা 5 এবং 5 এ-এর বিধানগুলি প্রযোজ্য হবে। একটি ব্যতিক্রম উপ-ধারা দ্বারা খোদাই করা হয়েছে।
\end{judgequote}

\begin{caselawref}
\bn{ডাব্লু. বি. এল. আর আইন, 1955-এর ধারা 3এ-এর 305 (2), যেখানে কোনও অকৃষি ভাড়াটিয়া বা আন্ডার-টেন্যান্ট কোনও জমির খাস দখল রাখে, সেক্ষেত্রে সে জমি রায়িয়াত হিসাবে ধরে রাখার অধিকারী। রায়িয়াতের দ্বারা জমি হস্তান্তরযোগ্যতা সম্পর্কিত বিধানও রয়েছে। যদি সেই জমিতে চাষাবাদ করা না হয়, তবে শ্রেণিবিন্যাসের জন্য একটি পরিবর্তন প্রয়োজন। 18। ডাব্লুবিএলআর আইনের ধারা 4-সি সম্পর্কিত জমি, 1955, কেবল মেমো নং-এর ভিত্তিতে সিদ্ধান্ত নেওয়া যায় না। ভি/তার অধিকার/775/15 তারিখটি উপ-জেলা জমি ও ভূমি সংস্কার অফিস, পুরূলিয়া দ্বারা জারি করা হয়েছিল। জমির মালিকদের দ্বারা প্রাপ্ত রাজস্ব অনুদানের কারণ হিসাবে নথিভুক্ত করা হয়নি। 1-বাগান্দিহর সর্বশ্রী চিরঞ্জীলাল (খনিজ) শিল্পগুলি এখনও ভাল রয়েছে, প্রাসঙ্গিক হবে কারণ স্থানান্তর, উত্তরাধিকার ইত্যাদির কারণে হাত বদল হতে পারে। এর সাথে সম্পর্কিত আইনি বিষয়গুলি। প্রথমত, বাগনডিহর 1 নং-সার্বশ্রী চিরঞ্জিলাল (খনি) ইন্ডাস্ট্রিজ তার অবস্থান পরিবর্তন করেছে কিনা এবং কেবলমাত্র 2016 সালের বিধি 61-এর সুবিধা পাওয়ার জন্য প্রয়োজনীয় তথ্য নিশ্চিত করতে পারবে না। 1-1998 সালে বাগান্দির সর্বশ্রী চিরঞ্জীলাল (খনিজ) শিল্পগুলি তখনও ভাল ছিল, যখন আবেদনটি দায়ের করা হয়েছিল, তখন কেন্দ্রীয় সরকারের অনুমোদনের প্রয়োজন ছিল। তবে, ডাব্লু. বি. এম. ডি. টি. সি. এল-কে একটি পক্ষ হিসাবে অন্তর্ভুক্ত করা হয়নি, যদিও এটি সর্বদা বাগান্দিহর 1 নং-সর্বেশ্রী চিরঞ্জিলাল (খনি) শিল্পের দ্বারা করা দাবির বিরোধিতা করে আসছিল, আমরা মনে করি না যে এই নির্দেশের আলোকে আমাদের এই বিষয়গুলি পরীক্ষা করার দরকার নেই।
\end{caselawref}

\noindent Appearances:


\begin{judgequote}
\bn{এটি বলার পরে, এটি আবেদনকারীদের অবস্থান-পশ্চিমবঙ্গ রাজ্য, যে তারা প্রশ্নবিদ্ধ জমির 20.87 একরের মালিক এবং এই রাজ্যের পক্ষে, তাদের যুক্তি প্রয়োগের সময়, সুনি ইজারা কার্যকর করার ক্ষেত্রে কোনও দ্বিধা নেই। এটি বর্ণিত অবস্থান হওয়ায়, যা আমাদের সামনেও যুক্ত করা হয়েছে, উক্ত অঞ্চলের জন্য খনি ইজারা দেওয়ার ক্ষেত্রে কোনও দ্বন্দ্ব থাকা উচিত নয়। তবে আমরা বলেছিলাম যে, 2015 সালের মাইনিং অ্যাক্ট নং 1-এর জন্যাকোব ডিপার্টমেন্টের এই সংশোধনী আইনের বিধানগুলি যথাযথভাবে পরীক্ষা করা হয়নি। সংশোধনী আইন, 2015 দ্বারা করা সংশোধনীগুলি উক্ত মামলায় সিদ্ধান্তের বিষয় ছিল না এবং আদালত ও কর্তৃপক্ষ কর্তৃক রায়টির অনুপাত নির্ভর করে প্রদত্ত রায়ে ব্যাখ্যা ও প্রয়োগ করা আইনি বিধানগুলির উপর নির্ভর করে। যখন আইনটি কোনও সংশোধনীর মাধ্যমে পরিবর্তিত হয়, তখন সংশোধিত আইনি বিধানগুলি হবে বিবেচনার বিষয়, ব্যাখ্যা করা এবং প্রয়োগ করা হয়। এবং বর্ণিত কারণগুলির জন্য, আমরা বর্তমান আপিলটি আংশিকভাবে অনুমোদন করি এবং রায়েরায়টি এই নির্দেশের সাথে বাতিল করে দিতে হবে যে পশ্চিমবঙ্গ সরকার আইডিএর <1 একর জমির জন্য একটি নির্দিষ্ট জায়গার জমি কার্যকর করবে, যা মিনিস্ট্রিজের পক্ষে থাকবে এবং 2030 সালের মধ্যে চিরঞ্জীবাগের দাবি হিসাবে বিবেচিত হবে। বর্তমান মামলার তথ্যে, খরচ সম্পর্কে কোনও আদেশ থাকবে না। দ্বারা প্রস্তুত করা শিরোনামঃ আংশিকভাবে অনুমোদিত আপিল। অঙ্কিত জ্ঞান 18 (2013) 9 এস. সি. সি 725।
\end{judgequote}

\noindent Respondent in Person.


\noindent JUDGMENT / ORDER OF THE SUPREME COURT


\noindent JUDGMENT


\noindent SANJIV KHANNA, J.


\begin{judgequote}
This appeal, by way of special leave, takes exception to the judgment
of the division bench of the High Court of Calcutta, whereby the intra-court
appeal preferred by the State of West Bengal and Others in F.M.A. No. 1458
\end{judgequote}


% --- page 6 ---

\noindent 282 SUPREME COURT REPORTS [2023] 12 S.C.R.


\begin{judgequote}
of 2017 with CAN No. 6596 of 2017 has been dismissed with the direction to
the Appellant No. 2 -- Joint Secretary, Department of Industries, Commerce
and Enterprises, West Bengal or any authorised offi cer to execute a mining
lease in favour of the Respondent No. 2 -- Dinesh Agarwal, sole proprietor
of Respondent No. 1 - M/s. Chiranjilal (Mineral) Industries of Bagandih.
\end{judgequote}

\begin{judgequote}
2. The facts are rather chequered, albeit are required to be noticed
in detail. On 07.08.1985, West Bengal Mineral Development and Trading
Corporation Limited 1 had fi led an application for grant of long term mining
lease for Dolomite, Limestone and Quartzite at the plots in Mouza -
Khariduara, Kumari and Boch. An application was also fi led by WBMDTCL
for grant of long term mining lease for Iron Ore, Manganese and Fireclay
at the plots in Mouza - Khariduara, Kumari, Boch and Kangametya. Grant
Order dated 07.04.1986 was issued in favour of WBMDTCL by the Assistant
Secretary, Commerce and Industries Department, Mines Branch, West
Bengal.
\end{judgequote}

\begin{judgequote}
2.1 On 06.03.1998, Respondent No. 1 - M/s. Chiranjilal (Mineral)
Industries of Bagandih. had fi led an application before the Mining Offi cer-
in-charge, Purulia Zone, Directorate of Mines and Minerals, West Bengal,
for the grant of a mining lease for the purpose of extracting Dolomite at
Mouza - Khariduara, Kumari and Boch, in 76 acres of land.
\end{judgequote}

\begin{judgequote}
2.2 The Respondent No. 1 - M/s. Chiranjilal (Mineral) Industries of
Bagandih fi led Writ Petition No. 7808 (W) of 2001 before the High Court
of Calcutta, seeking disposal of their application for grant of mining lease.
The High Court vide order dated 13.06.2001, directed the State authorities to
dispose of the application of Respondent No. 1 - M/s. Chiranjilal (Mineral)
Industries of Bagandih at an early date and in accordance with law.
\end{judgequote}

\begin{judgequote}
2.3. The Joint Secretary, Commerce and Industries Department, West
Bengal, vide order dated 13.03.2003, rejected the application of Respondent
No. 1 - M/s. Chiranjilal (Mineral) Industries of Bagandih, on the ground of
non-availability of land in view of the previous application of WBMDTCL.
By another order dated 26.03.2003, the Joint Secretary, Commerce and
Industries Department, West Bengal reiterated that the mining application
\end{judgequote}

\noindent 1 For Short,' WBMDTCL'.



% --- page 7 ---

\begin{judgequote}
STATE OF WEST BENGAL v. M/S. CHIRANJILAL (MINERAL) 283
INDUSTRIES OF BAGANDIH [SANJIV KHANNA, J.]
\end{judgequote}

\begin{judgequote}
of Respondent No. 1 - M/s. Chiranjilal (Mineral) Industries of Bagandih
overlaps with the area applied for in the previous application by WBMDTCL.
The application of the Respondent No. 1 - M/s. Chiranjilal (Mineral)
Industries of Bagandih was accordingly rejected.
\end{judgequote}

\begin{judgequote}
2.4 Aggrieved, the Respondent No. 1 - M/s. Chiranjilal (Mineral)
Industries of Bagandih had fi led Writ Petition No. 7505 (W) of 2003 in the
High Court of Calcutta challenging the orders passed by the Joint Secretary,
Commerce and Industries Department, West Bengal, dated 13.03.2003 and
26.03.2003. During the pendency of the said Writ Petition, the Joint Secretary,
Commerce and Industries Department, West Bengal, reviewed the aforesaid
orders and passed a fresh order dated 13.10.2006 forapportionment of land
between WBMDTCL and the Respondent No. 1 - M/s. Chiranjilal (Mineral)
Industries of Bagandih. This order states that two hearings were held on
24.05.2006 and 19.06.2006 to review the matter, and thereupon at the hearing
dated 19.06.2006, in the presence of the representatives of WBMDTCL and
the Respondent No. 1 - M/s. Chiranjilal (Mineral) Industries of Bagandih,
it was agreed that Respondent No. 1 - M/s. Chiranjilal (Mineral) Industries
of Bagandih will be granted the whole of the mining area of 76 acres, and
the lease for the rest of the area will be granted in favour of WBMDTCL.
No other reason has been stated and indicated in the said order. Thus, the
orders dated 13.03.2003 and 26.03.2003 rejecting the application of the
Respondent No. 1 - M/s. Chiranjilal (Mineral) Industries of Bagandih were
recalled. Consequently, the Letter of Intent dated 26.10.2006 was issued in
favour of the Respondent No. 1 - M/s. Chiranjilal (Mineral) Industries of
Bagandih for an area of 76 acres of land subject to fulfi lling/submission of
various documents, including approval of the Mining Plan duly approved
by the Chief Mining Offi cer, Asansol and Clearance Certifi cate from the
Ministry of Environment and Forests, Government of India.
\end{judgequote}

\begin{judgequote}
2.5 However, the order dated 13.10.2006 was cancelled or revoked
vide order dated 03.12.2010 by the Joint Secretary, Commerce and Industries
Department, Mines Branch, West Bengal, inter alia, recording that this
order was passed without ascertaining the exact position of the land and
in ignorance of the fact that the rejection orders dated 13.03.2003 and
26.03.2003 had already been challenged before the High Court in Writ
Petition No. 7505 (W) of 2003. The authorities had not ascertained the status
\end{judgequote}


% --- page 8 ---

\noindent 284 SUPREME COURT REPORTS [2023] 12 S.C.R.


\begin{judgequote}
of the case. The order of cancellation or revocation dated 03.12.2010 was
not challenged by the respondents.
\end{judgequote}

\begin{judgequote}
2.6. This order dated 03.12.2010 was also not brought to the notice
of the High Court, when the Writ Petition No. 7505 (W) of 2003 was
disposed of ex-parte vide order dated 25.03.2014 by relying upon the
supplementary affi davit fi led by the Respondent No. 1 - M/s. Chiranjilal
(Mineral) Industries of Bagandih, which had referred to the recalled order
dated 13.10.2006. This order of the High Court states that a decision as to
whether a lease or licence to be granted in favour of the Respondent No.
1 - M/s. Chiranjilal (Mineral) Industries of Bagandih shall be taken within
a period of eight weeks and Respondent No. 1 - M/s. Chiranjilal (Mineral)
Industries of Bagandih would be accordingly informed. It was made clear
that the decision as to the grant will be on the basis of the law and the rules
applicable at the time of consideration.
\end{judgequote}

\begin{judgequote}
2.7. By the order dated 09.07.2014 passed by the Joint Secretary,
Commerce and Industries Department, West Bengal, the application
fi led by the Respondent No. 1 - M/s. Chiranjilal (Mineral) Industries of
Bagandih was rejected inter alia relying upon the earlier application fi led
by WBMDTCL. Signifi cantly, this order mentions that the two rejection
orders dated 13.03.2003 and 26.03.2003 were recalled by the Joint Secretary
vide his order dated 13.10.2006. This order also refers to the factum that the
Grant Order dated 07.04.1986 to WBMDTCL forIron Ore, Manganese and
Fireclay in the plots in question had been revoked and the application for
Long-Term Mining Lease fi led by WBMDTCL for Dolomite and Limestone
was rejected by a common order dated 24.09.2009. The order dated
24.09.2009 has not been placed on record, though it is necessary to ascertain
and know the reasons for cancellation and rejection in favour of WBMDTCL.
WBMDTCL had applied earlier in point of time, and is a government of
West Bengal undertaking. The order dated 09.07.2014 does indicate that the
cancellation and rejection against WBMDTCL had something to do with
the Respondent No. 1 - M/s. Chiranjilal (Mineral) Industries of Bagandih,
and possibly the order dated 13.10.2006 in favour of the Respondent No.
1 - M/s. Chiranjilal (Mineral) Industries of Bagandih. This is refl ected from
the reason given in the order dated 09.07.2014, which states that since the
recall order dated 13.10.2006 was cancelled or revoked vide order dated
\end{judgequote}


% --- page 9 ---

\begin{judgequote}
STATE OF WEST BENGAL v. M/S. CHIRANJILAL (MINERAL) 285
INDUSTRIES OF BAGANDIH [SANJIV KHANNA, J.]
\end{judgequote}

\begin{judgequote}
03.12.2010, the rejection orders dated 13.03.2003 and 26.03.2003 were still
valid and the application for mining lease dated 07.08.1985 for Dolomite
and Limestone by WBMDTCL still subsists. Thereupon, reference in the
order dated 09.07.2014 is made to sub-section (2) to Section 11 2 of the Mines
\end{judgequote}

\begin{judgequote}
2 11. Preferential right of certain persons . - (1) Where a reconnaissance permit
or prospecting licence has been granted in respect of any land, the permit holder or the
licensee shall have a preferential right for obtaining a prospecting licence or mining lease,
as the case may be, in respect of that land over any other person:
Provided that the State Government is satisfi ed that the permit holder or the licensee, as the
case may be, -
(a) has undertaken reconnaissance operations or prospecting operations, as the case may be,
to establish mineral resources in such land;
(b) has not committed any breach of the terms and conditions of the reconnaissance permit
or the prospecting licence;
(c) has not become ineligible under the provision of this Act; and
(d) has not failed to apply for grant of prospecting licence or mining lease, as the case may
be, within three months after the expiry of reconnaissance permit or prospecting licence, as
the case may be, or within such further period as may be extended by the said Government.
(2) Subject to the provisions of sub-section (1),where the State Government has not
notifi ed in the Offi cial Gazette the area for grant of reconnaissance permit or prospecting
licence or mining lease, as the case may be, and two or more persons have applied for a
reconnaissance permit, prospecting licence or a mining lease in respect of any land in such
area, the applicant whose application was received earlier, shall have a preferential right to
be considered for grant of reconnaissance permit, prospecting licence or mining lease, as
the case may be, over the applicant whose application was received later:
Provided that where an area is available for grant of reconnaissance permit, prospecting
licence or mining lease, as the case may be, and the State Government has invited
applications by notifi cation in the Offi cial Gazette for grant of such permit, licence or
lease, all the applications received during the period specifi ed in such notifi cation and the
applications which had been received prior to the publication of such notifi cation in respect
of the lands within such area and had not been disposed of , shall be deemed to have been
received on the same day for the purposes of assigning priority under this subsection.
Provided further that where any such applications are received on the same day, the State
Government, after taking into consideration the matters specifi ed in sub-section (3), may
grant the reconnaissance permit, prospecting licence or mining lease, as the case may be, to
such one of the applicants as it may deem fi t.
(3) The matters referred to in sub-section (2) are the following:-
(a) any special knowledge of, or experience in, reconnaissance operations, prospecting
operations or mining operations, as the case may be, possessed by the applicant;
(b) the fi nancial resources of the applicant;
(c) the nature and quality of the technical staff employed or to be employed by the applicant;
(d) the investment which the applicant proposes to make in the mines and in the industry
based on the minerals;
(e) such other matters as may be prescribed.
\end{judgequote}


% --- page 10 ---

\noindent 286 SUPREME COURT REPORTS [2023] 12 S.C.R.


\begin{judgequote}
and Minerals (Development and Regulation) Act, 1957 3 , which states that
in cases where the State Government has not notifi ed in the Offi cial Gazette
an area for grant of reconnaissance permit, prospecting licence for mining
lease, and two or more persons had applied for the permit, licence or mining
lease, the person whose application received earlier in point of time shall
have preferential right for grant of permit, licence or lease over the person
whose application was received later. The order states that WBMDTCL is
very much interested in mining Dolomite and Limestone in the area and has
confi rmed the said fact in writing vide letter dated 05.06.2014.
\end{judgequote}

\begin{judgequote}
2.8. The Respondent No. 1 - M/s. Chiranjilal (Mineral) Industries
of Bagandih challenged the order dated 09.07.2014 passed by the Joint
Secretary, Commerce and Industries Department, West Bengal in Writ
Petition No. 21358 (W) of 2014 before the High Court of Calcutta. This
petition was disposed of vide order dated 10.09.2014 observing that the
Joint Secretary, who had passed the order dated 09.07.2014 had failed to
exercise jurisdiction vested in him as the applications fi led by WBMDTCL
had been rejected vide common order dated 24.09.2009 and were therefore
not pending. Direction was issued by the High Court to grant a long term
lease in respect of 76 acres of land to the Respondent No. 1 - M/s. Chiranjilal
(Mineral) Industries of Bagandih by observing that the respondent had a
Rayati status and that the remaining land can be given to WBMDTCL. It may
be relevant to note here that this order records that the fi les relating to the
\end{judgequote}

\begin{enumerate}[resume]
\item Subject to the provisions of sub-section(1), where the State Government notifi es in the
Offi cial Gazette an area for grant of reconnaissance permit, prospecting licence or mining
lease, as the case may be, all the applications received during the period as specifi ed in such
notifi cation, which shall not be less than thirty days, shall be considered simultaneously
as if all such applications have been received on the same day and the State Government,
after taking into consideration the matters specifi ed in sub-section(3), may grant the
reconnaissance permit, prospecting licence or mining lease, as the case may be, to such one
of the applicants as it may deem fi t.
(5) Notwithstanding anything contained in sub-section (2), but subject to the provisions of
sub-section (1), the State Government may, for any special reasons to be recorded, grant
a reconnaissance permit, prospecting licence or a mining lease, as the case may be, to
an applicant whose application was received later in preference to an applicant whose
application was received earlier:
Provided that in respect of minerals specifi ed in the First Schedule, prior approval of the
Central Government shall be obtained before passing any order under this sub-section.
3 For short, `MMDR Act, 1957'
\end{enumerate}


% --- page 11 ---

\begin{judgequote}
STATE OF WEST BENGAL v. M/S. CHIRANJILAL (MINERAL) 287
INDUSTRIES OF BAGANDIH [SANJIV KHANNA, J.]
\end{judgequote}

\begin{judgequote}
application of WBMDTCL were untraceable. WBMDTCL was not made a
party to the said writ petition. Notably, the application fi led by WBMDTCL,
being earlier in point of time in terms of the applicable rules was to be given
preference, whereas the application fi led by the Respondent No. 1 - M/s.
Chiranjilal (Mineral) Industries of Bagandih was rejected vide orders dated
13.03.2003 and 26.03.2003. However, the rejection orders were recalled vide
order dated 13.10.2006 and the Letter of Intent dated 26.10.2006 was issued
in favour of the Respondent No. 1 - M/s. Chiranjilal (Mineral) Industries
of Bagandih. Subsequently, the Grant Order dated 13.10.2006 in favour of
the Respondent No. 1 - M/s. Chiranjilal (Mineral) Industries of Bagandih
was cancelled and recalled vide order dated 13.12.2010. This order dated
13.12.2010 was never challenged and has attained fi nality. It is during the
period between the order dated 13.10.2006 and the order dated 13.12.2010
that the request/application of WBMDTCL was rejected and the mining
lease cancelled vide order dated 24.09.2009.
\end{judgequote}

\begin{judgequote}
2.9. On 10.02.2015, vide notifi cation No. S.O. 423 (E), Dolomite
was notifi ed as a minor mineral, and accordingly henceforth, fell under the
legislative and administrative jurisdiction of the State Government.
\end{judgequote}

\begin{judgequote}
2.10. A Grant Order dated 16.07.2015 was issued by the Deputy
Secretary, Commerce and Industries Department, West Bengal for Dolomite
mining in favour of Respondent No. 1 - M/s. Chiranjilal (Mineral) Industries
of Bagandih in respect of 76 acres of land, subject to certain conditions,
including the requirement to submit consent letters of owners of the land
in question ( Raiyats ) before the execution of the lease deed, or a condition
to this eff ect would be incorporated in the draft lease. Another stipulation
mentioned therein is the need for permission under Section 14-Y 4 of the West
\end{judgequote}

\begin{judgequote}
4 14-Y. Limitation on future acquisition of land by a raiyat. ---If at any time, after
the commencement of the provisions of this Chapter, the total area of land owned by
a raiyat exceeds the ceiling area applicable to him under Section 14-M, on account of
transfer, inheritance or otherwise, the area of land which is in excess of the ceiling area
shall vest in the State and all the provisions of this Chapter relating to ceiling area shall
apply to such land:
Provided that a person intending to establish a tea garden, mill, factory or workshop,
livestock breeding farm, poultry farm, or dairy, or township in accordance with the
provisions of the West Bengal Town and Country (Planning and Development) Act, 1979,
may, with the previous permission, in writing, of the State Government and on such terms
\end{judgequote}


% --- page 12 ---

\noindent 288 SUPREME COURT REPORTS [2023] 12 S.C.R.


\begin{judgequote}
Bengal Land Reforms Act, 1955 5 for holding the required land and furnishing
of Conversion Certifi cate for plots of land from the appropriate authority in
terms of Section 4-C 6 of the WBLR Act, 1955. It is also stipulated that the
Grant Order and the subsequent execution of the lease deed are subject to
the No Objection Certifi cate to be obtained from the Central Government
since Dolomite was a major mineral at the time of the order dated 10.09.2014
passed by the High Court.
\end{judgequote}

\begin{judgequote}
2.11. Aggrieved by the conditions and the requirements stipulated in
the Grant Order dated 16.07.2015, the Respondent No. 1 - M/s. Chiranjilal
(Mineral) Industries of Bagandih fi led two Contempt Petitions in W.P.
\end{judgequote}

\begin{judgequote}
and conditions and in such manner as the State Government may by rules prescribe, acquire
and hold land in excess of the ceiling area applicable to him under Section 14-M:
Provided further that if such person, having been permitted by the State Government,
does not utilise within two years of the date of such permission such land for the purpose
for which he has been so permitted by the State Government to acquire and hold it, then, all
the provisions of this Chapter relating to ceiling area shall apply to the area of land which
is held in excess of the ceiling area applicable to him under Section 14-M.
Explanation .---For the purpose of this section, ``person'' includes an individual,
a fi rm, a company, an institution, or an association or body of individuals, whether
incorporated or not.
5 For short, `WBLR Act, 1955'.'
6 4-C. Permission for change of area, character or use of land. ---(1) A raiyat holding
any land may apply to the Collector for change of area or character of such land or for
conversion of the same for any purpose other than the purpose for which it was settled or
was being previously used or for alteration in the mode of use of such land.
(2) On receipt of such application, the Collector may, after making such inquiry
as may be prescribed and after giving the applicant or the persons interested in such land
or aff ected in any way an opportunity of being heard, by order in writing either reject the
application or direct such change, conversion or alteration, as the case may be, on such
terms and conditions as may be prescribed.
(3) Every order under sub-section (2) directing change, conversion or alteration shall
specify the date from which such change, conversion or alteration shall take eff ect.
(4) A copy of the order passed by the Collector directing change, conversion or
alteration, if any, under sub-section (2), or in an appeal therefrom shall he forwarded to
the Revenue Offi cer referred to in Section 50 or Section 51, as the case may be, and such
Revenue Offi cer shall incorporate in the record-of-rights changes eff ected by such order
and revise the record-of-rights in accordance with such order.
(5) If the Collector is satisfi ed that any land is being convened for any purpose other
than the purpose for which it was settled or was being previously held, or attempts are being
made to eff ect alteration in the mode of use of such land or change of the area or character
of such land, he may, by order, restrain the raiyat from such Act.
\end{judgequote}


% --- page 13 ---

\begin{judgequote}
STATE OF WEST BENGAL v. M/S. CHIRANJILAL (MINERAL) 289
INDUSTRIES OF BAGANDIH [SANJIV KHANNA, J.]
\end{judgequote}

\begin{judgequote}
21358 (W) of 2014. These contempt petitions were disposed of, inter alia ,
observing that the Respondent No. 1 - M/s. Chiranjilal (Mineral) Industries
of Bagandih was required to fulfi l the conditions, including furnishing of the
Conversion Certifi cate under Section 4-C of the WBLR Act, 1955 and No
Objection Certifi cate from the Government of India. The court, therefore,
found that there was no wilful, or contumacious violation of the order dated
10.09.2014. However, liberty was granted to the Respondent No. 1 - M/s.
Chiranjilal (Mineral) Industries of Bagandih to question the Grant Order
dated 16.07.2015.
\end{judgequote}

\begin{judgequote}
2.12. The Respondent No. 1 - M/s. Chiranjilal (Mineral) Industries of
Bagandih thereupon preferred Writ Petition No. 20309 (W) of 2016 before
the High Court of Calcutta. However, WBMDTCL was not a party to this
writ petition. In the meanwhile, a clarifi cation was sought by the Deputy
Secretary, Commerce and Industries Department, West Bengal and vide
clarifi cation dated 26.08.2016 issued by the Government of India, Ministry
of Mines, it was clarifi ed that even prior to 10.02.2015, Dolomite was a
Non-Scheduled major mineral, for which prior approval of the Central
Government was not required under sub-section (1) to Section 5 of the
MMDR Act, 1957.
\end{judgequote}

\begin{judgequote}
2.13.This Writ Petition No. 20309 (W) of 2016 vide judgment and
order dated 12.04.2017 has been allowed inter alia observing that Dolomite
had become a minor mineral with eff ect from 10.02.2015 and hence prior
approval of the Central Government is not required under Section 5(1) of
the MMDR Act, 1957. On the question of requirements under Section 14-Y
and 4-C of the WBLR Act, 1955, it is observed that the land in question is
recorded as `Dungri' as per information provided by the Deputy District
Land and Land Reforms Offi cer, Purulia vide Memo No. V/RTI/775/15
dated 06.03.2017 and that the land classifi ed as `Dungri' is only used for
the purpose of mining lease and thus, there is no need for a conversion
certifi cate under Section 4-C of the WBLR, Act, 1955. The clarifi cation dated
07.04.2016 was issued by the Additional District Magistrate and District
Land and Land Reforms Offi cer, Purulia, stating that the Respondent No.
1 - M/s. Chiranjilal (Mineral) Industries of Bagandih had procured a No
Objection Certifi cate in respect of the major portion of Raiyati land from
diff erent owners and that the State Government itself was the owner of
\end{judgequote}


% --- page 14 ---

\noindent 290 SUPREME COURT REPORTS [2023] 12 S.C.R.


\begin{judgequote}
20.87 acres of land, thus Section 14-Y of the WBLR Act, 1955 would not be
applicable as the Respondent No. 1 - M/s. Chiranjilal (Mineral) Industries
of Bagandih has not acquired land in excess ceiling limit prescribed under
Section 14-M of the WBLR Act, 1955.
\end{judgequote}

\begin{judgequote}
2.14. This judgment was challenged by the State of West Bengal in
an intra-court appeal being F.M.A. No. 1458 of 2017 with CAN No. 6596
of 2017 which has been dismissed vide the impugned judgment dated
04.10.2018. Agreeing with the fi ndings recorded by the Single Judge,
the division bench has held that the provisions of the West Bengal Minor
Minerals Concession Rules, 2016 7 will not be applicable as the Respondent
No. 1 - M/s. Chiranjilal (Mineral) Industries of Bagandih had made the
application in March 1998, and more so as the Joint Secretary, Government
of West Bengal had passed the order dated 13.10.2006 to grant mining
lease. The High Court's direction given in Writ Petition No. 21358 (W) of
2014 vide judgment dated 10.09.2014 are prior to the enforcement of the
Concession Rules, 2016.
\end{judgequote}

\begin{judgequote}
3. We have heard the learned Senior Advocate appearing for the State
of West Bengal and the Respondent No. 2 -- Dinesh Agarwal, who has
appeared in-person. They have also submitted their written submissions.
\end{judgequote}

\begin{judgequote}
4. We begin our discussion by first referring to Rule 61 of the
Concession Rules, 2016, which reads as under:
\end{judgequote}

\begin{judgequote}
``61.Decleration of ineligibility of the pending minor mineral
applications for mining lease including the applications of
reclassifi ed major minerals.- All applications for mining lease of
minor minerals including the reclassifi ed minor minerals vide SO No-
423 (E) dated 12 th February,2015 received prior to the giving-eff ect
to this rules irrespective of its duration of pendency shall become
ineligible.
\end{judgequote}

\begin{judgequote}
Provided that if the applicant has been issued a Grant Order or Letter
of Intent (LoI) or any other Government Order requiring the alteration
of applicant's position then his mining lease application may be
considered after due compliance of the all the necessary conditions
\end{judgequote}

\noindent 7 For short, `Concession Rules, 2016'.



% --- page 15 ---

\begin{judgequote}
STATE OF WEST BENGAL v. M/S. CHIRANJILAL (MINERAL) 291
INDUSTRIES OF BAGANDIH [SANJIV KHANNA, J.]
\end{judgequote}

\begin{judgequote}
5. An almost corresponding amendment was made to the MMDR Act,
1957 by incorporating Section 10-A vide Mines and Minerals (Development
and Regulation) Amendment Act, 2015 8 , which reads as under:
\end{judgequote}

\begin{judgequote}
10-A. Rights of existing concession holders and applicants.--- (1)
All applications received prior to the date of commencement of the
Mines and Minerals (Development and Regulation) Amendment Act,
2015, shall become ineligible.
\end{judgequote}

\begin{enumerate}[resume]
\item Without prejudice to sub-section (1), the following shall remain
eligible on and from the date of commencement of the Mines and
Minerals (Development and Regulation) Amendment Act, 2015---
\end{enumerate}

\noindent (a) applications received under Section 11-A of this Act;


\begin{judgequote}
(b) where before the commencement of the Mines and Minerals
(Development and Regulation) Amendment Act, 2015 a reconnaissance
permit or prospecting licence has been granted in respect of any land
for any mineral, the permit holder or the licensee shall have a right
for obtaining a prospecting licence followed by a mining lease, or a
mining lease, as the case may be, in respect of that mineral in that
land, if the State Government is satisfi ed that the permit holder or the
licensee, as the case may be,---
\end{judgequote}

\begin{judgequote}
(i) has undertaken reconnaissance operations or prospecting operations,
as the case may be, to establish the existence of mineral contents in
such land in accordance with such parameters as may be prescribed
by the Central Government;
\end{judgequote}

\begin{judgequote}
(ii) has not committed any breach of the terms and conditions of the
reconnaissance permit or the prospecting licence;
\end{judgequote}

\noindent (iii) has not become ineligible under the provisions of this Act; and


\begin{judgequote}
(iv) has not failed to apply for grant of prospecting licence or mining
lease, as the case may be, within a period of three months after the
expiry of reconnaissance permit or prospecting licence, as the case
may be, or within such further period not exceeding six months as
may be extended by the State Government;
\end{judgequote}

\noindent 8 For short, `Amendment Act, 2015'.



% --- page 16 ---

\noindent 292 SUPREME COURT REPORTS [2023] 12 S.C.R.


\begin{judgequote}
(c) where the Central Government has communicated previous
approval as required under sub-section (1) of Section 5 for grant of
a mining lease, or if a letter of intent (by whatever name called) has
been issued by the State Government to grant a mining lease, before
the commencement of the Mines and Minerals (Development and
Regulation) Amendment Act, 2015, the mining lease shall be granted
subject to fulfi lment of the conditions of the previous approval or
of the letter of intent within a period of two years from the date of
commencement of the said Act:
\end{judgequote}

\begin{judgequote}
Provided that in respect of any mineral specifi ed in the First Schedule,
no prospecting licence or mining lease shall be granted under clause
(b) of this sub-section except with the previous approval of the Central
Government.
\end{judgequote}

\begin{judgequote}
6. Rule 61 of the Concession Rules, 2016 states that all applications
for mining lease of minor minerals including reclassifi ed minor minerals
vide S.O. No. 423 (E) dated 12.02.2015 received prior to giving eff ect
to the Concession Rules, 2016 9 , irrespective of its duration of pendency
shall become ineligible. In other words, these applications are not to be
considered. The proviso makes an exception and states that if an applicant,
who had made an application prior to 29.07.2016, had been issued a Grant
Order or a Letter of Intent, or any other order requiring alteration of the
applicant's position, his application for mining lease may be considered
after due compliance of all necessary conditions. The question is whether
the respondents' case is covered by the exception in terms of the proviso
to Rule 61 of the Concession Rules,2016. We have already referred to the
reasoning given by the division bench of the High Court dealing with the
Concession Rules, 2016, and would like to quote the fi ndings which hold
that the proviso would not be applicable to the facts of the present case.
These observations read:
\end{judgequote}

\begin{judgequote}
``25. \textbackslash{}ldots\{\}.Neither such recent policy nor can the provisions of the West
Bengal Minor Minerals Concession Rules, 2016 can apply to the
application of the writ petitioners made in March, 1998 and more so
as the order of the Joint Secretary to grant lease is dated 13 th October,
\end{judgequote}

\noindent 9 The Concession Rules, 2016 came into eff ect on 29.07.2016



% --- page 17 ---

\begin{judgequote}
STATE OF WEST BENGAL v. M/S. CHIRANJILAL (MINERAL) 293
INDUSTRIES OF BAGANDIH [SANJIV KHANNA, J.]
\end{judgequote}

\begin{judgequote}
2006 and that of this Court directing grant of long term lease is dated
10 th September, 2014 are prior to such policy and prior to the said
Rules came into operation. It further appears that necessary mining
plan taking into account the environmental aspect has been submitted
by the writ petitioners and the appellant/State has raised no grievance
in respect thereof.
\end{judgequote}

\begin{judgequote}
7. The policy referred to in the aforesaid paragraph is in terms of the
letter dated 02.02.2018 issued by the Principal Secretary, State of West
Bengal, wherein it is specifi ed that obtaining a Conversion Certifi cate is a
mandatory condition for the purpose of a mining lease. Reference in the
impugned judgment to the order dated 13.10.2006, or for that matter, the
Letter of Intent dated 26.10.2006 is inconsequential as the said orders were
recalled and revoked on 03.12.2010. The orders did not survive and continue
to operate thereafter. Writ Petition No. 7505 (W) of 2003 was disposed of
ex-parte, without noticing that the order dated 13.10.2006 had been recalled
or cancelled, albeit the judgment had directed that the application for grant of
lease would be considered in accordance with law and the rules applicable at
the time of consideration. The order dated 03.12.2010 was never challenged
by the Respondent No. 1 - M/s. Chiranjilal (Mineral) Industries of Bagandih
and has attained fi nality. At best, the case of the Respondent No. 1 - M/s.
Chiranjilal (Mineral) Industries of Bagandih is that the application dated
06.03.1998 should be considered in accordance with law.
\end{judgequote}

\begin{judgequote}
8. The Respondent No. 1 - M/s. Chiranjilal (Mineral) Industries of
Bagandih has relied upon judgment of this Court in Bhushan Power and
Steel Limited v. S.L. Seal, Additional Secretary (Steel and Mines), State
of Odisha and Others 10 . In the said case, the predecessor-in-interest of the
petitioner therein had made an application for grant of lease before the State
of Odisha for mining of Iron Ore in an area measuring 1250 acres. The
application was in view of the proposal to set up a steel plant in the district
of Sambalpur, Odisha. The rejection for the grant of the mining lease to M/s.
Bhushan Power and Steel Limited was challenged in a Writ Petition in the
High Court, which was dismissed, but the appeal preferred before this Court
was allowed vide judgment dated 14.03.2012 in Bhushan Power and Steel
\end{judgequote}

\noindent 10 (2017) 2 SCC 125.



% --- page 18 ---

\noindent 294 SUPREME COURT REPORTS [2023] 12 S.C.R.


\begin{judgequote}
Limited and Others v. State of Orissa and Another 11 , setting aside the order
of the State Government dated 09.02.2016, with the following directions:
``41 . In the light of the above, the High Court erred in holding that it
could not interfere with the decision of the State Government calling
upon the appellants to sign a fresh MoU with the Government, during
subsistence of the earlier MoU. Since the State Government has already
made allotments in favour of others in relaxation of the Mineral
Concession Rules, 1960, under Rule 59(2) thereof, no cogent ground
had been made out on behalf of the State to deny the said privilege
to the appellants as well. Accordingly, we allow the appeal and set
aside the judgment and order of the High Court of Orissa and also
the decision of the State Government dated 9-2-2006, rejecting the
appellants' claim for grant of mining lease.
42 . During the course of hearing, we have been informed that Thakurani
Block A has large reserves of iron ore, in which the appellants can
also be accommodated. We, accordingly, direct the State of Orissa to
take appropriate steps to act in terms of the MoU dated 15-5-2002, as
also its earlier commitments to recommend the case of the appellants
to the Central Government for grant of adequate iron ore reserves to
meet the requirements of the appellants in their steel plant at Lapanga.
\end{judgequote}

\begin{judgequote}
9. The State of Odisha thereafter fi led an application for review of the
judgment in Bhushan Power and Steel Limited and Others v. State of Orissa
and Another 12 (supra)which was rejected vide order dated 11.09.2012.
\end{judgequote}

\begin{judgequote}
10. Alleging non-compliance and in-action of the judgment dated
14.03.2012, a contempt petition was fi led by M/s Bhushan Power and Steel
Limited. The contempt petition was contested by the State of Odisha on
several grounds, including that the judgment dated 14.03.2012 is incapable
of enforcement, for which reliance was placed on a subsequent judgment of
this Court in Sandur Manganese and Iron Ores Ltd. v. State of Karnataka 13 .
This stand did not fi nd favour with this Court and the offi cers of the State
Government were found to be in contempt of the judgment dated 14.03.2012
\end{judgequote}

\begin{judgequote}
11 (2012) 4 SCC 246.
12 (2012) 4 SCC 246.
13 (2010) 13 SCC 1.
\end{judgequote}


% --- page 19 ---

\begin{judgequote}
STATE OF WEST BENGAL v. M/S. CHIRANJILAL (MINERAL) 295
INDUSTRIES OF BAGANDIH [SANJIV KHANNA, J.]
\end{judgequote}

\begin{judgequote}
vide judgment dated 22.04.2014 in Bhushan Power and Steel Limited v.
Rajesh Verma 14 . Under these circumstances, the judgment dated 22.04.2014
had given one more opportunity to the State Government to send requisite
recommendation to the Central Government inter alia observing that this
Court cannot lose sight of the fact that there is a judgment inter se the parties,
which has become fi nal. Accordingly, the contention that the judgment of
this Court in Sandur Manganese and Iron Ores Limited (supra) will not
undo the directions given in the judgment dated 14.03.2012 was rejected.
The relevant observations in the judgment dated 22.04.2014 read as under:
\end{judgequote}

\begin{judgequote}
`` 21. We cannot lose sight of the fact that there is a judgment, inter
partes, which has become fi nal. Even when the civil appeal was being
heard, certain other parties claiming their interest in these very lands
had moved intervention applications which were dismissed. At that
time also it was mentioned that there are 195 applicants. However,
notwithstanding the same, this Court issued fi rm directions to the
State Government to recommend the case of the petitioners for mining
lease in both the areas. In view of such categorical and unambiguous
directions given in the judgment which has attained fi nality, merely
because another judgment has been delivered by this Court in Sandur
Manganese case , cannot be a ground to undo the directions contained
in the judgment dated 14-3-2012. Insofar as law laid down in Sandur
Manganese is concerned, that may be applied and followed by the State
Government in respect of other applications which are still pending.
However, that cannot be pressed into service qua the petitioner whose
rights have been crystallised by the judgment rendered in its favour.
It cannot be reopened, that too at the stage of implementation of the
said judgment.
\end{judgequote}

\begin{judgequote}
22. \textbackslash{}ldots\{\}. Once we hold that the respondents are bound to implement the
direction contained in the judgment dated 14-3-2012, insofar as the
State Government is concerned, it is obliged to comply therewith and
such matters, along with other relevant considerations, can be left to
the wisdom of the Central Government while taking a decision on the
recommendation of the State Government.
\end{judgequote}

\noindent 14 (2014) 5 SCC 551.



% --- page 20 ---

\noindent 296 SUPREME COURT REPORTS [2023] 12 S.C.R.


\noindent xx xx xx


\begin{judgequote}
24 . \textbackslash{}ldots\{\}. However, we are giving one fi nal opportunity to them to purge
the contempt by transmitting requisite recommendations to the Central
Government. It would be for the Central Government to consider the
said recommendations on its own merits and in accordance with law.
In case the recommendation is sent within one month from the date
of copy of receipt of this order, we propose not to take any further
action and the respondents/contemnors shall stand discharged from
this contempt petition. However, in case the respondents do not purge
in the manner mentioned above, it would be open to the petitioners
to point out the same to this Court by moving appropriate application
and in that event the contemnors shall be proceeded against.
\end{judgequote}

\noindent (emphasis supplied)


\begin{judgequote}
11. Consequent to the directions dated 22.04.2014, the State Government
had sent the requisite recommendation to the Central Government for grant
of mining lease of the area in question. The Central Government, however,
took the stand that having regard to the amendments in the MMDR Act,
1957, vide the Amendment Act, 2015 introducing Section 10-A, the request
made by M/s Bhushan Power and Steel Limited stands invalidated. In
view of the aforesaid stand, the Central Government had written letters to
the State Government, with a copy sent to M/s Bhushan Power and Steel
Limited. In the letter dated 13.05.2015, the Central Government had stated
that the proposal for according the prior approval for grant of mineral
concession was ineligible in terms of sub-section (1) to Section 10-A of
the MMDR Act, 1957 and, therefore, should be treated as closed. However,
the State Government might ascertain whether the proposal was safe from
ineligibility under Section 10-A of the MMDR Act, 1957 and thereupon
the State Government could take action accordingly. Similar view was also
expressed by the Central Government in the letter dated 29.05.2015 therein.
Consequent to these communications, the State government vide letter dated
09.07.2015 had informed M/s Bhushan Power and Steel Limited that their
applications for grant of mining lease had become ineligible as per sub-
section (1) to Section 10-A of the MMDR Act, 1957.
\end{judgequote}


% --- page 21 ---

\begin{judgequote}
STATE OF WEST BENGAL v. M/S. CHIRANJILAL (MINERAL) 297
INDUSTRIES OF BAGANDIH [SANJIV KHANNA, J.]
\end{judgequote}

\begin{judgequote}
12. This Court in M/s Bhushan Steel and Power Limited 15 (supra),
specifi cally examined the contention whether in the facts of the said case,
clause (c) to sub-section (2) to Section 10-A of the MMDR Act, 1957 could
be invoked in view of the contention raised by M/s Bhushan Steel and
Power Limited that the Letter of Intent was issued by the State Government
for grant of mining lease and, therefore, their application stands protected.
The submission was that the recommendation dated 24.05.2014, given by
the State Government should be treated as a Letter of Intent by ``whatever
name called'', as it signifi es the intention to grant mining lease insofar as
the State Government is concerned. It was also argued that under the new
regime contained under Section 10-A of the MMDR Act, 1957, approval of
the Central Government was not even required and the State Government
could have proceeded further and granted the lease.
\end{judgequote}

\begin{judgequote}
13. The aforesaid arguments did not fi nd favour of this Court in the
case of M/s Bhushan Steel and Power Limited (supra) in spite of the
earlier judgment of this Court dated 14.03.2012 and the order passed in the
contempt petition dated 22.04.2014 with the observations therein that there
was failure of the State Government to comply with the directions. This
Court rejected the submissions in M/s Bhushan Steel and Power Limited
(supra) and held as under:
\end{judgequote}

\begin{judgequote}
``17. Undoubtedly, as per sub-section (1) of Section 10-A, all
applications received prior to coming into force of the Amendment
Act, 2015, become ineligible. Reason for interpreting such a provision
is not far to seek. Before the passing of the Amendment Act, 2015, it
was the Central Government which had the ultimate control over the
grant of licences insofar as mining of major minerals is concerned.
As per the procedure then existing, the State Government could
recommend the application submitted by any applicant for grant of
mining lease to the Central Government and the Central Government
was given the power to grant or refuse to grant the approval. Thus,
``previous approval'' from the Central Government was essential for
grant of lease, without which the State Government could not enter
\end{judgequote}

\noindent 15 (2012) 4 SCC 246.



% --- page 22 ---

\noindent 298 SUPREME COURT REPORTS [2023] 12 S.C.R.


\begin{judgequote}
into any such lease agreement with the applicant. Shortcomings of
this procedure were noticed by this Court in its judgment rendered
in Centre for Public Interest Litigation v. Union of India [(2012) 3 SCC
1] (for short `` CPIL case '') and also in Natural Resources Allocation,
In re, Special Reference No. 1 of 2012 [Natural Resources Allocation,
In re, Special Reference No. 1 of 2012, (2012) 10 SCC 1] . In these
judgments, this Court expressed that allocation of natural resources
should normally be by auction. Judgment in CPIL case had a direct
relevance to the grant of mineral concessions as the Government found
that it was resulting in multipurpose litigation which was becoming
counterproductive. Mining Ordinance, 2015 was passed on 12-1-
2015 which was ultimately replaced when Parliament enacted the
Amendment Act, 2015.
\end{judgequote}

\begin{judgequote}
18. The exhaustive Statement of Objects and Reasons reveals that
the extensive amendment in the Act were eff ected after extensive
consultations and intensive scrutiny by the Standing Committee on
Coal and Steel, who gave their Report in May 2013. As is evident from
the Statement that diffi culties were experienced because the existing
Act does not permit the auctioning of mineral concessions. It was
observed that with auctioning of mineral concessions, transparency in
allocation will improve; the Government will get an increased share of
the value of mineral resources; and that it will alleviate the procedural
delay, which in turn would check slowdown which adversely aff ected
the growth of mining sector.
\end{judgequote}

\begin{judgequote}
19. The Amendment Act, 2015, as is evident from the objects, aims
at: ( i ) eliminating discretion; ( ii ) improving transparency in the
allocation of mineral resources; ( iii ) simplifying procedures; ( iv )
eliminating delay on administration, so as to enable expeditious and
optimum development of the mineral resources of the country; ( v )
obtaining for the Government an enhanced share of the value of the
mineral resources; and ( vi ) attracting private investment and the latest
technology.
\end{judgequote}

\begin{judgequote}
20. The Amendment Act, 2015 ushered in the amendment of Sections 3,
4, 4-A, 5, 6, 13, 15, 21 and First Schedule; substitution of new sections
for Sections 8, 11 and 13; and, insertion of new Sections 8-A, 9-B,
\end{judgequote}


% --- page 23 ---

\begin{judgequote}
STATE OF WEST BENGAL v. M/S. CHIRANJILAL (MINERAL) 299
INDUSTRIES OF BAGANDIH [SANJIV KHANNA, J.]
\end{judgequote}

\begin{judgequote}
9-C, 10-A, 10-C, 11-B, 11-C, 12-A, 15-A, 17-A, 20-A, 30-B, 30-C
and Fourth Schedule.
\end{judgequote}

\begin{judgequote}
21. These amendments brought in vogue: ( i ) auction to be the sole
method of allotment; ( ii ) extension of tenure of existing lease from the
date of their last renewal to 31-3-2030 (in the case of captive mines)
and till 31-3-2020 (for the merchant miners) or till the completion
of renewal already granted, if any, or a period of 50 years from the
date of grant of such lease; ( iii ) establishment of District Mineral
Foundation for safeguarding interest of persons aff ected by mining
related activities; ( iv ) setting up of a National Mineral Exploration Trust
created out of contributions from the mining lease-holders, in order
to have a dedicated fund for encouraging exploration and investment;
( v ) removal of the provisions requiring `` previous approval '' from
the Central Government for grant of mineral concessions in case of
important minerals like iron ore, bauxite, manganese, etc. thereby
making the process simpler and quicker; ( vi ) introduction of stringent
penal provisions to check illegal mining prescribing higher penalties
up to Rs 5 lakhs per hectare and imprisonment up to 5 years; and ( vii )
further empowering the State Government to set up Special Courts
for trial of off ences under the Act.
\end{judgequote}

\begin{judgequote}
14. Thus, the object and purpose of the Amendment Act, 2015 is to
ensure that allocation of mineral resources is done through auctioning. This
is the reason why sub-section (1) to Section 10-A of the MMDR Act, 1957
mandates that all applications received prior to 12.01.2015 shall become
ineligible. The exceptions or the saving clause applies to three kinds of
situations specifi ed in sub-section (2) to Section 10-A of the MMDR Act.
1957. The fi rst category is where an application has been received under
Section 11-A of the MMDR Act,1957. The second category is where
a reconnaissance permit or a prospecting licence has been granted the
permit holder or the licensee has the right to obtain a prospecting licence
followed by a mining lease and the State Government is satisfi ed that the
permit holder or the licensee has complied with the requirements specifi ed
in sub-clauses (i) to (iv) of clause (b) of sub-section (2) to Section 10-A of
the MMDR Act, 1957. The reason for protecting this class of cases is on
account of the fact that they had altered their position by spending money
\end{judgequote}


% --- page 24 ---

\noindent 300 SUPREME COURT REPORTS [2023] 12 S.C.R.


\begin{judgequote}
on reconnaissance operations or prospecting operations. Accordingly, the
principle of legitimate expectation is applied. The third category is where
the Central Government had already communicated their previous approval
or the State Government had issue Letter of Intent for grant of mining lease
before coming into force of the Amendment Act 2015. The raison dêtre,
it is observed therein, is that certain rights had accrued to these applicants
inasmuch as all necessary procedures and formalities had been complied
with and only formal lease remains to be executed.
\end{judgequote}

\begin{judgequote}
15. Delving on the question of whether the letter for approval dated
22.05.2014 granted by the State Government can be treated as a Letter of
Intent predicated on the words by whatever name, which expression, it was
submitted, should be given a broad interpretation in view of the words `by
whatever name called' was examined in-depth and in detail. Reference was
made to the legal dictionary for the meaning of the term `Letter of Intent'
as a preliminary understanding between the parties who intend to make
a contract or join together for further action. Reference was also made to
decisions of this Court in Rishi Kiran Logistics Private Limited v. Board
of Trustees of Kandla Port Trust and Others 16 and Rajasthan Cooperative
Dairy Federation Limited v. Maha Laxmi Mingrate Marketing Service
Private Limited and Others 17 However, the said contention was rejected
inter alia holding as under:
\end{judgequote}

\begin{judgequote}
`` 26. Applying the aforesaid meaning, can it be said that Letter dated
24-5-2014 of the State Government would constitute a letter of intent?
We are afraid, answer has to be in the negative. Reason is simple. As
mentioned above, in order to enable the State Government to enter into
any lease agreement/contract with the prospecting licensee, `` previous
approval '' of the Central Government was essential. Unless such
approval came, the State Government could not communicate to the
prospecting licensee/lessee its intention to enter into any contract as the
prerequisite prior approval would be lacking. Therefore, no promise
could be held by the State Government to any applicant showing its
intention to enter into a contract in the future. Position would have
\end{judgequote}

\begin{judgequote}
16 (2015) 13 SCC 233.
17 (1996) 10 SCC 405.
\end{judgequote}


% --- page 25 ---

\begin{judgequote}
STATE OF WEST BENGAL v. M/S. CHIRANJILAL (MINERAL) 301
INDUSTRIES OF BAGANDIH [SANJIV KHANNA, J.]
\end{judgequote}

\begin{judgequote}
been diff erent had Letter dated 24-5-2014 been issued after receiving
previous approval of the Central Government. However, that is not so.
This letter to the Central Government was only recommendatory in
nature and ultimate decision rested with the Central Government. It is a
diff erent thing if the Central Government refuses to give its approval on
any extraneous reasons or mala fi des or does not take into consideration
relevant factors/material while rejecting the application, which may
form a diff erent cause of action and may become a reason to challenge
the action of the Central Government rejecting the application on the
grounds that are available in law to seek judicial review of such an
action. However, we are not dealing with that situation in the instant
case. Our discussion is confi ned to the plea raised before us viz.
whether Letter dated 24-5-2014 can be termed as ``letter of intent''.
For the reasons stated above, we are of the view that it was not a letter
of intent. The application of the petitioner, therefore, would not be
covered by clause ( c ) of Section 10-A of the Act.
\end{judgequote}

\begin{judgequote}
27. We are conscious of the fact that the petitioner herein had originally
succeeded in the appeal inasmuch as judgment dated 14-3-2012 was
rendered giving direction to the State Government to recommend
the case of the petitioner, in terms of the MoU entered into between
the parties, to the Central Government. This was not done and the
decision was reiterated in orders dated 22-4-2014 passed in Bhushan
Power and Steel Ltd. v. Rajesh Verma [. It is possible that had the
State Government acted promptly and sent the recommendations
earlier, the Central Government might have accorded its approval.
However, whether it could have done so or not would be in the realm
of conjectures. Insofar as the Central Government is concerned,
no direction was ever given by this Court. On the contrary, it was
categorically observed in the order dated 22-4-2014 in Bhushan
Power and Steel Ltd. v. Rajesh Verma that it would be for the Central
Government to consider the recommendations of the State Government
on its own merits and in accordance with law. If that has not been done
by the Central Government, it cannot be the subject-matter of present
contempt petition.
\end{judgequote}

\begin{judgequote}
16. The aforesaid judgment is relevant for our purpose, though in the
present case, post notifi cation No. S.O. 423(E) dated 10.02.2015, Dolomite
\end{judgequote}


% --- page 26 ---

\noindent 302 SUPREME COURT REPORTS [2023] 12 S.C.R.


\begin{judgequote}
was notifi ed as a minor mineral and hence, the approval of the Central
Government was not required for the reason that the Grant Order dated
16.07.2015 was hedged with pre-conditions, including the requirement
to submit consent letters of the owners of the land in question ( Raiyats )
before the execution of the lease deed, or there was to be a stipulation that a
condition to this eff ect would be incorporated in the draft lease. Therefore, in
our opinion, the Grant Order dated 16.07.2015 is provisional, and is subject
to fulfi lment of the conditions therein. This is clear from the terms of the
Grant Order dated 16.07.2015, which are reproduced below:
\end{judgequote}

\begin{judgequote}
``
xx xx xx
\end{judgequote}

\begin{judgequote}
(a) You have to furnish a Draft Mining Lease Deed in the model form
K us prescribed in the Mineral Concession Rules, 1960, as amended
upto date (1 rule 31 of MC Rules, 1960),
\end{judgequote}

\begin{judgequote}
(b) The Draft Mining Lease Deed should be prepared in durable papers
neatly and suffi cient space should be kept in between two lines in order
to permit, if necessary, correction therein,
\end{judgequote}

\begin{judgequote}
(c) The Deed of Lease, after execution, shall be registered by you
at your own cost and no mining operation should be started before
registration of the Deed,
\end{judgequote}

\begin{judgequote}
(d) You shall have to furnish the approved Mining Plan, if not submitted
rules 22(4) and 22A of MC Rules, 1960,
\end{judgequote}

\begin{judgequote}
(e) You shall have to furnish the Environment Clearance (EC), if not
submitted from the M1EF Environment of Protection Act, 1986,
\end{judgequote}

\begin{judgequote}
(f) You shall have to furnish Consent to Establish and Consent to
Operate from the WBPCB before execution of Deed of Lease [Section
25 and 26 of Water Act ,1974 and Section 21 of Air Act, 1981],
\end{judgequote}

\begin{judgequote}
(g) You shall have to raise annually a minimum quantity of minerals
as stipulated in the approved Mine Plan [rules 22A and 45(ia) of MC
Rules, 1960],
\end{judgequote}

\begin{judgequote}
(h) You shall have to deposit Rs. 10,000/ (Rupees ten thousand) only
as Security for due observance of the terms and conditions of the lease,
\end{judgequote}


% --- page 27 ---

\begin{judgequote}
STATE OF WEST BENGAL v. M/S. CHIRANJILAL (MINERAL) 303
INDUSTRIES OF BAGANDIH [SANJIV KHANNA, J.]
\end{judgequote}

\begin{judgequote}
under appropriate Head of Account which shall be refundable to you
after expiry of the period of Lease, unless the whole or a part of it is
withheld or forfeited by the Government for any default on you part
including default in payment of amount due to the Government [rule
32 of MC Rules, 1960],
\end{judgequote}

\begin{judgequote}
(i) You shall have to submit consent letter(s) of the owner(s) of the
land under consideration before execution of the Lease Deed (Consent
of the Raiyats) or a condition to that eff ect should be incorporated in
the Draft Deed (rule 22(3)(i)(1t)),
\end{judgequote}

\begin{judgequote}
(j) You shall have to furnish the N.O.C., of the Forest Authority in
proper format in case the applied area falls in the forest area as notifi ed
by the Appropriate Authority, alongwith the Draft Lease Deed or
a condition to that eff ect should be incorporated in the Draft Deed
[Section 2 of Forest Conservation Act, 1980],
\end{judgequote}

\begin{judgequote}
(k) For actual operation of quarrying or digging, ten (10) yards clear
margin shall be kept from the outer boundary of the adjacent 1 plot
or plots and maintain throughout the operation and you shall have to
give a written undertaking to that eff ect or corporate a condition in
the Draft Lease Deed,
\end{judgequote}

\begin{judgequote}
(l) You shall have to incorporate all the conditions as mentioned in
the M.C. Rules, 1960 in the Draft Lease Deed,
\end{judgequote}

\begin{judgequote}
(m) You shall have to furnish up to date Royalty Clearance, Income Tac
Clearance and VAT Clearance certifi cates before execution of the Deed,
\end{judgequote}

\begin{judgequote}
(n) You shall have to submit, along with the Draft Deed, a Geo-
Reference Map duly vetted by the DL\&LRO and DMM, West Bengal,
if not submitted,
\end{judgequote}

\begin{judgequote}
(o) You shall have obtained the permission under Section 14Y of
WBLR Act, 1955 for holding the required land,
\end{judgequote}

\begin{judgequote}
(p) You have to furnish the Conversion Certifi cate for plots of land
from the Appropriate Authority (Section 4C of WBLR Act, 1955),
\end{judgequote}

\begin{judgequote}
(q) You have to furnish the current Land Availability Report (LAR)
from the Appropriate Authority.
\end{judgequote}


% --- page 28 ---

\noindent 304 SUPREME COURT REPORTS [2023] 12 S.C.R.


\begin{judgequote}
(r) In the event of non-execution of the deed within the stipulated
period on compliance with the above mentioned conditions the order
sanctioning the lease shall be liable to be revoked,
\end{judgequote}

\begin{judgequote}
(s) You shall have to comply with all the statutory requirements before
presenting the Deed of Lease of execution to this Department,
\end{judgequote}

\begin{judgequote}
(t) This Grant Order and subsequent execution of Lease Deed are
subject to the No Objection Certifi cate (NOC) to be obtained by this
Department form the Govt. of India since the applicant prayed for
mining lease on the ground that the Letter of Intent (Lol) was issued
for the mineral Dolomite which was a major mineral at the time of
order dated 10.09.2014 of the Hon'ble High Court.
\end{judgequote}

\noindent xx xx xx''


\begin{judgequote}
17. Raiyat land is to be used for cultivation, etc., and not for mining.
Once the mining activity is undertaken, the Raiyats will not be able to use
the land. In terms of sub-section (10) to Section 2 of the WBLR Act, 1955,
a Raiyat means a person or an institution holding land for any purpose
whatsoever. However, the rights of Raiyat in respect of the land in terms of
sub-section (2A) to Section 4 of the WBLR Act, 1955 does not permit any
other person to quarry sand from his holding, dig or use, or permit any person
to dig or use, earth or clay of his holding for the manufacture of bricks or
tiles except with previous permission in writing of the State Government.
In case of breach of the condition, the prescribed authority may, after giving
notice and opportunity to a Raiyat to show cause, can levy a monetary
penalty. Further, on an order being passed, the land shall vest in the State
free from all encumbrances. Section 4-B of the WBLR Act, 1955 stipulates
that every Raiyat holding any land shall maintain and preserve such land in
a manner that the area is not diminished or its character is not changed or
the land is not converted for any purpose other than the purpose for which
it was settled or previously held except with the previous permission of the
Collector in writing. Equally signifi cant for our purpose is Section 3A of
the WBLR Act, 1955, which states that the rights and interests of all non-
agricultural tenants and under-tenants shall vest in the State free from all
encumbrances and provisions of Section 5 and 5A of the West Bengal Estates
Acquisition Act, 1953 shall apply. An exception is carved out by sub-section
\end{judgequote}


% --- page 29 ---

\begin{judgequote}
STATE OF WEST BENGAL v. M/S. CHIRANJILAL (MINERAL) 305
INDUSTRIES OF BAGANDIH [SANJIV KHANNA, J.]
\end{judgequote}

\begin{enumerate}[resume]
\item to Section 3A of the WBLR Act, 1955, where a non-agricultural tenant
or under-tenant is holding khas possession of any land, in which case he is
entitled to retain the land as Raiyat . There are also provisions relating to the
transferability of land by the Raiyat. If cultivation was not being undertaken
at the land in question, the classifi cation requires a change.
\end{enumerate}

\begin{judgequote}
18. The controversy relating to Section 4-C of the WBLR Act, 1955,
cannot simply be decided on the basis of Memo No. V/RTI/775/15 dated
06.03.2017 issued by the Deputy District Land and Land Reforms Offi cer,
Purulia, that as per the revenue records the land was recorded as ` Dungri '.
The reason is that Raiyat land is not for mining. Thus, a contradiction
arises, as the grant of Raiyat land and the classifi cation of the same land as
` Dungri ' is contradictory.
\end{judgequote}

\begin{judgequote}
19. Further, whether the consent letter of the owners of the land in
question ( Raiyats ) obtained by the Respondent No. 1 - M/s. Chiranjilal
(Mineral) Industries of Bagandih still hold good, would be relevant as there
could be a change of hands on account of transfer, inheritance, etc. Connected
with this are the legal issues. First, whether the Respondent No. 1 - M/s.
Chiranjilal (Mineral) Industries of Bagandih had altered its position post the
issue of the Grant Order dated 16.07.2015, but before enforcement of the
Concession Rules, 2016, to get the benefi t of Rule 61 of the Concessions
Rules, 2016? It is necessary to ascertain the facts and then alone one can
adjudicate and decide the question whether the Respondent No. 1 - M/s.
Chiranjilal (Mineral) Industries of Bagandih is entitled to the benefi t of
the proviso to Rule 61 of the Concession Rules, 2016. This has not been
verifi ed and ascertained. An issue would arise on whether the application
fi led by the Respondent No. 1 - M/s. Chiranjilal (Mineral) Industries of
Bagandih way back in 1998 would still hold good as at the time, when the
application was fi led, approval of the Central Government was required.
Another diffi culty is that WBMDTCL has not been impleaded as a party,
though it was always contesting the claim made by the Respondent No.
1 - M/s. Chiranjilal (Mineral) Industries of Bagandih. On the question of
cancellation or rejection of the application made by WBMDTCL, we have
made observations supra. However, we need not examine these issues in
light of the order and directions we are issuing. Further, we feel that the
remand order should not be passed at this distinct point of time.
\end{judgequote}


% --- page 30 ---

\noindent 306 SUPREME COURT REPORTS [2023] 12 S.C.R.


\begin{judgequote}
20. Having said so, it is the stand of the appellants -- State of West
Bengal, that they are owners of 20.87 acres of the land in question and to
this extent, they have no diffi culty in executing the mining lease. This being
the stated stand, which has also been affi rmed before us, there should be
no diffi culty in granting of mining lease for the said area to the Respondent
No. 1 - M/s. Chiranjilal (Mineral) Industries of Bagandih.
\end{judgequote}

\begin{judgequote}
21. During the course of arguments before us, reference was made by
the appellants to the provisions of the WBLR Act, 1955 and the judgment
of this Court in Thressiamma Jacob and Others v. Geologist, Department
of Mining and Geology and Others 18 . We have not examined the said
aspects which are left open and not adjudicated upon. However, we deem
it appropriate to observe that the judgment of this Court in Thressiamma
Jacob and Others (supra) is prior to the enforcement of the Amendment
Act, 2015 and the Concession Rules, 2016. The amendments made by
the Amendment Act, 2015 were not subject matter of decision in the said
case and would have to be considered by the courts and the authorities as
a judgment's binding ratio depends upon the legal provisions considered,
interpreted and applied in a given judgment. When the law changes by
an amendment in the legislation, the amended legal provisions have to be
considered, interpreted and applied.
\end{judgequote}

\begin{judgequote}
Accordingly, and for the reasons stated, we partly allow the present
appeal and set aside the impugned judgment with a direction that the
government of West Bengal will execute a mining lease for 20.87 acres of
land in favour of the Respondent No. 1 - M/s. Chiranjilal (Mineral) Industries
of Bagandih. The Writ Petition No. 20309 (W) of 2016 will be treated as
allowed to the extent as indicated above. The claim of the Respondent No.
1 - M/s. Chiranjilal (Mineral) Industries of Bagandih towards the balance
area for the grant of mining lease will be treated as rejected and dismissed.
In the facts of the present case, there will be no order as to costs.
\end{judgequote}

\begin{judgequote}
Headnotes prepared by: Appeal partly allowed.
Ankit Gyan
\end{judgequote}

\noindent 18 (2013) 9 SCC 725.



\end{document}
